\documentclass[11pt]{article}
\usepackage[T1]{fontenc}
\usepackage{lmodern}
\usepackage[margin=1in]{geometry}
\usepackage{amsmath,amssymb,amsthm,mathtools}
\usepackage{microtype,booktabs,array,longtable}
\usepackage[dvipsnames]{xcolor}
\usepackage{hyperref}
\hypersetup{colorlinks=true,linkcolor=MidnightBlue,citecolor=MidnightBlue,urlcolor=MidnightBlue,
 pdftitle={Relative Discriminant Charges and Signed Source Response in Theta Heat Flow},
 pdfauthor={Drafted for Edward Baker}}
\usepackage{bookmark}
\newtheorem{theorem}{Theorem}[section]
\newtheorem{lemma}[theorem]{Lemma}
\newtheorem{proposition}[theorem]{Proposition}
\newtheorem{corollary}[theorem]{Corollary}
\theoremstyle{definition}
\newtheorem{definition}[theorem]{Definition}
\newtheorem{example}[theorem]{Example}
\theoremstyle{remark}
\newtheorem{remark}[theorem]{Remark}
\numberwithin{equation}{section}
\setlength{\emergencystretch}{1.5em}
\allowdisplaybreaks[1]
\title{Relative Discriminant Charges and Signed Source Response\\in Theta Heat Flow}
\author{Drafted for Edward Baker}
\date{21 October 2026}
\begin{document}
\maketitle
\begin{abstract}
We develop a finite-cluster framework for studying multiple zeros of the
backward heat evolution of the Riemann xi function. The local degree of
the holomorphic jet $(H,H_z)$ equals the winding of the cluster
discriminant; an event of spatial multiplicity $m$ has charge
$\binom m2$. For a proper two-root cluster, the discriminant satisfies a
regular affine equation $\Delta'=8+A\Delta$. A complex-time phase bound
reduces its charge to the signs at two real time caps. We construct
explicit complete-source windows at every hypothetical finite-height
double, derive a signed discriminant response with a quadratic remainder,
and control event displacement along an affine heat homotopy. The
canonical theta tail has a positive Laplace-mixture representation that
strictly improves a matched pointwise-envelope enclosure, without yet
determining the required sign. For arbitrary finite clusters we prove
$D'=8T+BD$ and exhibit exact cubic and quartic obstructions to a direct
extension of the two-root cap rule. Existing interval certificates and
exact arithmetic checks provide supporting examples. The missing
theta-specific signed estimate, its height coverage, and the higher
multiplicity exclusion remain explicit open obligations. No endpoint
exclusion or RH-conditional simplicity theorem is claimed.
\end{abstract}
\clearpage
\begingroup
\small
\tableofcontents
\endgroup
\clearpage
\section{The endpoint problem and the complete source}
\label{sec:intro}

The purpose of this paper is to assemble the analytic results of the
relative-charge investigation into a single argument with explicit
hypotheses. The question motivating the investigation is whether the
unchanged theta source supplies enough signed information to exclude a
multiple zero at time zero, conditional on the Riemann hypothesis (RH).
The results below supply geometric windows, exact response identities,
and quantitative estimates to which such information could be applied.
They also identify limitations of several simpler substitutes.

\subsection{Normalization and analytic bounds}

Throughout the theta application, put
\begin{align}
 H(t,z)&=\int_0^\infty \Phi(u)e^{tu^2}\cos(zu)\,du,\label{i:source}\\
 \Phi(u)&=\sum_{n\geq1}\phi_n(u),\qquad \lambda_n=\pi n^2,\notag\\
 \phi_n(u)&=(2\lambda_n^2e^{9u}-3\lambda_ne^{5u})
                  e^{-\lambda_ne^{4u}}.\label{i:phi}
\end{align}
This is the standard normalization used in \cite{polymath}:
\begin{equation}
 H(0,z)=\frac18\xi\!\left(\frac12+\frac{iz}{2}\right),\qquad
 \xi(s)=\frac{s(s-1)}2\pi^{-s/2}\Gamma(s/2)\zeta(s).
 \label{i:xi}
\end{equation}
The Fourier representation \eqref{i:xi} and the classical real-zero
threshold theorem are the analytic number-theory inputs; none of the
local proofs below uses a zero table or RH unless it is explicitly stated.

For $u\geq0$ every summand in \eqref{i:phi} is positive, since
$2\lambda_ne^{4u}-3>0$. The series and its derivatives converge
uniformly on compact real $u$ intervals. For large $u$ the summands are
bounded by a polynomial in $n,e^u$ times $e^{-\pi n^2e^{4u}}$;
factoring out, for example, $e^{-(\pi/2)e^{4u}}$ leaves a uniformly
summable series. Consequently $\Phi$ and each of its fixed-order
derivatives decay faster than $e^{-cu^2-du}$ for any fixed real $c,d$.
Differentiation under the integral is therefore justified locally
uniformly in $(t,z)\in\mathbb C^2$. In particular,
\begin{equation}
 H_t=-H_{zz},\qquad
 H(t,-z)=H(t,z),\qquad
 H(\overline t,\overline z)=\overline{H(t,z)}.
 \label{i:heat}
\end{equation}
Here and below subscripts denote derivatives; the time parameter is
written as an argument to avoid confusing it with a derivative.

For $T\in\mathbb R$, $R\geq0$ and nonnegative integer $j$, define the
\emph{complete moment}
\begin{equation}
 M_j(T,R)=\int_0^\infty \Phi(u)u^j e^{Tu^2}\cosh(Ru)\,du.
 \label{i:moments}
\end{equation}
Since $|\cos(x+iy)|$ and $|\sin(x+iy)|$ are at most $\cosh y$,
\begin{equation}
 |\partial_t^a\partial_z^bH(t,z)|\leq M_{2a+b}(T,R)
 \quad(\operatorname{Re}t\leq T,\ |\operatorname{Im}z|\leq R).
 \label{i:jet-bound}
\end{equation}
These bounds are independent of $\operatorname{Re}z$. ``Complete'' means
that every source mode and the entire integration tail are included.
For a signed residual source $\psi$, the same assertion holds with
$|\psi|$ in place of $\Phi$.

\subsection{What RH supplies, and what it leaves open}

The de Bruijn--Newman theorem gives a finite real $\Lambda$ for which
all zeros of $H(t,\cdot)$ are real exactly when $t\geq\Lambda$.
In this normalization RH is equivalent to $\Lambda\leq0$;
Rodgers and Tao proved $\Lambda\geq0$ \cite{rodgers-tao}.
Thus under RH the zeros are real for all $t\geq0$. They are simple for
$t>0$: if a real zero had multiplicity $m\geq2$ at $t_0>0$, the local
Hermite splitting proved in Proposition~\ref{f:hermite}, applied at
$t=t_0-\varepsilon$, would give nonreal zeros for sufficiently small
$\varepsilon>0$. At least two Hermite roots are nonzero, and their
imaginary displacement has order $\sqrt{\varepsilon}$, whereas the
error is $O(\varepsilon)$. This contradicts reality at that time.
The argument does not rule out a multiple zero at $t=0$.

A common negative time at which \emph{all} theta zeros are real is
unavailable, even under RH, since it would contradict $\Lambda\geq0$.
This restriction does not prohibit different negative-time neighborhoods
for different endpoint heights. The quantitative windows below are
explicitly allowed to depend on the hypothetical event.

The local charge at an endpoint of multiplicity $m$ will be
$q=\binom m2$. It relates directly to the multiplicity excess studied
in the companion heat-motion investigation \cite{bw10,assessment07}.
Under RH let $\rho=\tfrac12+i\gamma$ range over distinct nontrivial
zeros of both ordinate signs, with multiplicity $m_\rho$. For every
finite symmetric cutoff $Y$, define
\begin{equation}
 K_Y=\sum_{|\gamma|\leq Y}\frac{m_\rho^2}{|\rho|^2},\qquad
 Z_Y=\sum_{|\gamma|\leq Y}\frac{m_\rho}{|\rho|^2}.
 \label{i:weighted}
\end{equation}
Termwise algebra gives
\begin{equation}
 K_Y-Z_Y=2\sum_{|\gamma|\leq Y}\frac{q_\rho}{|\rho|^2}.
 \label{i:bridge}
\end{equation}
When the two weighted sums converge, as in the companion framework,
the same identity holds without the cutoff. No convergence theorem is
needed for the local arguments in this paper. At the physical coordinate
$a=2\gamma$, the charge weight on the right is $8/(1+a^2)$.
Thus the local integer charge does not produce a positive uniform gap
for the weighted excess. An endpoint-local charge must also be
distinguished from the charge of a larger window, which might contain a
negative-time event.

\subsection{Degree, relative Chern number, and Hopf interpretation}

Let $\Omega\subset\mathbb C^2$ be a bounded domain with piecewise smooth
boundary, and suppose $G=(H,H_z)$ is nonzero on that boundary. Equip both
copies of $\mathbb C^2$ with the complex orientation, in the coordinate
order $(\operatorname{Re}t,\operatorname{Im}t,
\operatorname{Re}z,\operatorname{Im}z)$. The boundary charge is
\begin{equation}
 Q_\Omega=\deg\left(\frac{G}{\|G\|}:\partial\Omega\longrightarrow S^3\right).
 \label{i:degree}
\end{equation}
It is the relative top Chern number of the trivial complex rank-two
bundle with the specified nonvanishing boundary section: equivalently,
the total index of its interior zeros. For isolated holomorphic zeros
these indices are positive analytic intersection multiplicities.
Theorem~\ref{f:charge} supplies the discriminant computation of them.

For the product of a time rectangle and a spatial disk, rounding corners
makes the boundary an oriented $S^3$. Its usual Hopf interpretation
requires no additional charge formula. Let $\pi:S^3\to S^2$ be the Hopf
map, let $\omega$ have integral one on $S^2$, and choose a connection
form $\alpha$ with $d\alpha=\pi^*\omega$ and
$\int_{S^3}\alpha\wedge d\alpha=1$. For $g=G/\|G\|$ the Hopf invariant
of $\pi\circ g$ is
\[
 \int_{\partial\Omega}g^*\alpha\wedge d(g^*\alpha)
 =\int_{\partial\Omega}g^*(\alpha\wedge d\alpha)=\deg g.
\]
Here the integral formula defines the normalized Hopf invariant; its
independence of the chosen primitive follows by Stokes' theorem and
$H^1(S^3)=0$. It agrees with \eqref{i:degree}. The analytic estimates
in this paper use degree and discriminants directly.

Neither reality nor spatial parity supplies cancellation. Conjugating
both complex domain coordinates preserves real orientation, as does
conjugating both target coordinates. The change $z\mapsto-z$ and its
target action $(H,H_z)\mapsto(H,-H_z)$ are also orientation preserving.
The corresponding reflected events therefore have the same positive
local charge.

\subsection{Reading the results}

The first part proves proper cluster preparation, local charge, the
two-root cap theorem and explicit event-compatible windows. The second
part retains signs in source perturbations, proves quantitative event
continuation, and specializes the response to the actual canonical theta
tail. The many-root part keeps every root of a higher birth and shows
exactly where the two-root argument ceases to apply. Numerical evidence
is collected separately in Section~\ref{sec:support}. Appendix~\ref{app:reference}
records the holomorphic reference comparison and the boundary
interpolation contract used by earlier certificates.

All propositions are statements with the hypotheses displayed here.
In particular, a narrower source-response interval need not lie on one
side of zero, and a phase bound need not imply zero charge. This
manuscript consolidates internally reviewed research results; it has not
received independent specialist validation. It asserts no priority for
the standard complex-analysis and heat-root-motion ingredients.

\section{Finite clusters and their positive charge}\label{f:cluster-section}

The objects in this section are finite clusters, rather than a product
over all zeros of an entire function.  This distinction makes the
boundary assumptions explicit and retains the influence of the exterior
zeros through a holomorphic unit.  The construction consolidates
\cite{note04}.

\begin{definition}[Proper spatial cluster]\label{f:proper}
Let $D\subset\mathbb C$ be a bounded disk or rectangle, and let
$\Omega_R=\{z:|z-a|<R\}$.  Suppose $H$ is holomorphic on a neighborhood
of $\overline D\times\overline\Omega_R$ and
\[
 H(t,z)\ne0\qquad
 (t\in\overline D,\ |z-a|=R).
\]
The zeros of $H(t,\cdot)$ in $\Omega_R$, counted with multiplicity,
form a proper cluster.  Their number $m$ is independent of $t$.
A collar of radius $r<R$ means that all these zeros lie in
$\Omega_r=\{|z-a|<r\}$ throughout $\overline D$.
\end{definition}

The count is the integer-valued continuous function
$(2\pi i)^{-1}\oint H_z/H\,dz$.  In particular, two counts, one on each
of the circles of radii $r$ and $R$, certify a collar; an inner count
alone does not exclude additional roots in the annulus.

\begin{lemma}[Preparation without root labels]\label{f:preparation}\label{f:cluster}
A proper cluster has a monic polynomial
\[
 P(t,z)=z^m+p_1(t)z^{m-1}+\cdots+p_m(t)
\]
with coefficients holomorphic in $t$, and
$H(t,z)=U(t,z)P(t,z)$, where $U$ is holomorphic and nonvanishing on
the product.  The polynomial is uniquely determined by the cluster.
\end{lemma}

\begin{proof}
For $k\ge1$ the contour integrals
\begin{equation}\label{f:root-moments}
 S_k(t)=\frac1{2\pi i}\oint_{|\zeta-a|=R}
       \zeta^k\frac{H_z(t,\zeta)}{H(t,\zeta)}\,d\zeta
\end{equation}
are holomorphic and equal the sums of the $k$th powers of the roots,
including multiplicity.  Newton's identities determine the elementary
symmetric functions recursively from $S_1,\ldots,S_m$, hence construct
$P$ with holomorphic coefficients.  On each fixed time fiber $H/P$
has removable singularities at exactly the selected roots and has no
zeros in the disk.  Joint holomorphy follows, without choosing those
roots, from the Cauchy representation
\[
 U(t,z)=\frac1{2\pi i}\oint_{|\zeta-a|=R}
        \frac{H(t,\zeta)}{P(t,\zeta)(\zeta-z)}\,d\zeta,
       \qquad |z-a|<R.
\]
For each $t$ this is the removable quotient, and the integrand is
holomorphic in its parameters.  Near the spatial wall the quotient
itself supplies the extension.  Compactness and the strict boundary
condition allow the same construction on a slightly enlarged time
neighborhood.  Uniqueness follows from the multiset of roots.
\end{proof}

Define the finite discriminant by
\begin{equation}\label{f:discriminant-definition}
 \Delta_P(t)=\operatorname{Disc}_z P(t,z)
    =\prod_{i<j}(z_i(t)-z_j(t))^2.
\end{equation}
The symmetric expression is holomorphic even when the roots cannot be
labeled holomorphically.  We set $\Delta_P=1$ if $m\le1$.

\begin{theorem}[Boundary charge and discriminant winding]
\label{f:charge-winding}\label{f:charge}
In addition to properness, suppose that $(H,H_z)$ has no zero on
$\partial D\times\overline\Omega_R$.  Give complex two-space its
complex orientation and the product boundary its induced orientation.
The degree of
\[
 \frac{(H,H_z)}{\sqrt{|H|^2+|H_z|^2}}:
       \partial(D\times\Omega_R)\longrightarrow S^3
\]
is well defined and equals
\begin{equation}\label{f:winding-formula}
 Q_{D\times\Omega_R}
   =\frac1{2\pi i}\oint_{\partial D}
              \frac{\Delta_P'(t)}{\Delta_P(t)}\,dt.
\end{equation}
It is the sum of the local complex intersection multiplicities of
$(H,H_z)$, and is therefore a nonnegative integer.
\end{theorem}

\begin{proof}
The side wall contains no zero of $H$, and the remaining boundary
contains no joint zero by hypothesis.  The discriminant is nonzero
on $\partial D$.  It is consequently not identically zero and has
only finitely many zeros in $\overline D$.  Above each such time there
are only finitely many cluster roots, so the joint zeros are isolated.

Fix one such time $t_0$ and put
$\mathcal R=\mathbb C\{t-t_0\}$, the ring of convergent power series.
The module $\mathcal A=\mathcal R[z]/(P)$ is free of rank $m$ over
$\mathcal R$.  Multiplication by $P_z$ on $\mathcal A$ has determinant
\[
 \operatorname{Res}_z(P,P_z)
    =(-1)^{m(m-1)/2}\Delta_P.
\]
Indeed, at a time with distinct roots multiplication is diagonal in
the evaluation coordinates, with determinant $\prod_iP_z(z_i)$;
the identity then extends holomorphically.  The ring $\mathcal R$ is
a discrete valuation ring.  Elementary row and column reduction puts
any matrix with nonzero determinant into diagonal form with entries
units times powers $(t-t_0)^{e_j}$.  Thus the order of its determinant
is the sum of the $e_j$, which is also the length of its cokernel.
Applying this to multiplication by $P_z$ gives
\[
 \operatorname{ord}_{t_0}\Delta_P
       =\operatorname{length}_{\mathcal R}\mathcal A/(P_z).
\]
Small disjoint spatial disks around the distinct roots at $t_0$
split this quotient into its local factors.  In each factor, analytic
division by the monic local polynomial identifies its length with
\[
 \dim_{\mathbb C}
 \frac{\mathbb C\{t-t_0,z-b\}}{(P,P_z)}.
\]
Since $U$ is a unit, $(H,H_z)=(UP,UP_z+U_zP)$ and $(P,P_z)$
generate the same ideal.  These dimensions are therefore precisely
the local intersection multiplicities of the original jet.

For completeness, the topological interpretation uses the usual local
degree property of a holomorphic square map: at an isolated zero its
degree is the dimension of its local analytic quotient.  One way to
see the sign and multiplicity is to perturb the target to a nearby
regular value in an isolating neighborhood.  The local intersection
number counts the preimages with their multiplicities; at a regular
preimage the real Jacobian determinant is
$|\det_{\mathbb C}DF|^2>0$.  Conservation of the intersection number
and homotopy invariance of boundary degree give the asserted equality
at the original value.  Additivity over disjoint isolating neighborhoods
then identifies the product-boundary degree with the sum of the above
lengths.  The scalar argument principle sums
$\operatorname{ord}_{t_0}\Delta_P$ and proves
\eqref{f:winding-formula}.  Corners of a product rectangle may be
rounded within a boundary neighborhood free of joint zeros, without
changing either degree.
\end{proof}

The heat equation determines the local multiplicity in this formula.
Here and below the time normalization is $H_t=-H_{zz}$.

\begin{proposition}[Hermite birth and local charge]\label{f:hermite}
Suppose $H(t_0,b+w)=c_mw^m+O(w^{m+1})$, where $c_m\ne0$ and
$m\ge2$.  The event $(t_0,b)$ is an isolated zero of $(H,H_z)$.
For its local $m$-root cluster, writing $s=t-t_0$ and
$N=m(m-1)/2$, one has
\begin{equation}\label{f:hermite-discriminant}
 \Delta_P(t_0+s)=d_m s^N+O(s^{N+1}),\qquad
 d_m=2^N\prod_{i<j}(\lambda_i-\lambda_j)^2>0,
\end{equation}
where $\lambda_1,\ldots,\lambda_m$ are the roots of the probabilists'
Hermite polynomial $\operatorname{He}_m$.  The local jet charge is $N$.
\end{proposition}

\begin{proof}
The heat equation recursively determines the Taylor coefficients in
$s$ from the initial spatial Taylor series.  Terms arising from $w^\ell$
have weighted degree $\ell$ when $w$ has weight one and $s$ has weight
two.  The leading weighted polynomial is
\[
 c_m e^{-s\partial_w^2}w^m
   =c_m(2s)^{m/2}\operatorname{He}_m
                  \left(\frac{w}{\sqrt{2s}}\right).
\]
This is a polynomial identity, so the notation involving square roots
introduces no choice into the polynomial itself.

The polynomial $\operatorname{He}_m$ has $m$ distinct real roots.
Indeed, its Rodrigues formula and integration by parts imply
orthogonality to every polynomial of degree less than $m$ for the
positive weight $e^{-x^2/2}$.  If it had fewer than $m$ distinct real
roots of odd multiplicity, multiply it by the product of the linear
factors at just those roots.  The resulting real polynomial would
have a constant nonzero sign except at finitely many points, so its
weighted integral could not vanish.  Orthogonality contradicts this.
Degree $m$ now forces all $m$ roots to be real and simple.

Set $s=\varepsilon^2$ and $w=\varepsilon x$.  Taylor convergence gives,
uniformly for bounded $x$,
\[
 \varepsilon^{-m}H(t_0+\varepsilon^2,b+\varepsilon x)
 =c_m2^{m/2}\operatorname{He}_m(x/\sqrt2)+O(\varepsilon).
\]
The left side extends holomorphically to $\varepsilon=0$.
The simple-root implicit function theorem gives all $m$ cluster roots
as
$w_i=\varepsilon(\sqrt2\lambda_i+O(\varepsilon))$.
They are distinct for sufficiently small nonzero $\varepsilon$.
Their squared differences have product
$d_m\varepsilon^{2N}(1+O(\varepsilon))$.
The discriminant is holomorphic in $s=\varepsilon^2$, so its Taylor
series has only even powers of $\varepsilon$; this proves
\eqref{f:hermite-discriminant}.  It also proves local isolation of
the event.  Theorem~\ref{f:charge-winding} identifies its charge with
the order $N$.
\end{proof}

Consequently, on any product satisfying Theorem~\ref{f:charge-winding}
for a heat family,
\begin{equation}\label{f:positive-sum}
 Q_{D\times\Omega_R}
    =\sum_{(t,b):\,H(t,b)=H_z(t,b)=0}\binom{m(t,b)}2.
\end{equation}
Interior event charges do not cancel.  Reality and spatial parity do
not change this: simultaneous conjugation in two complex coordinates
has real orientation sign $+1$, and spatial sign reversal is complex
linear, also of sign $+1$.  These symmetries can pair equal charges,
but cannot force their sum to vanish.

\section{The two-root equation and the real-cap theorem}
\label{f:two-root-section}

Assume a proper cluster consists of two roots in a collar
$0<r<R$.  Write $H=UP$, $P=(z-z_1)(z-z_2)$, and
$\Delta=(z_1-z_2)^2$.  Set $q=U_z/U$ and, on the outer circle,
$f=H_z/H$.  The exterior unit $U$ is retained throughout; replacing
it by a constant would change the equation for the original cluster.

\begin{proposition}[Regular affine discriminant flow]\label{f:affine}
The holomorphic coefficient
\begin{equation}\label{f:A-contour}
 A(t)=\frac4{2\pi i}\oint_{|\zeta-a|=R}
        \frac{H_z(t,\zeta)}{H(t,\zeta)P(t,\zeta)}\,d\zeta
\end{equation}
is regular through collisions and satisfies
\begin{equation}\label{f:affine-equation}
                    \Delta'=8+A\Delta.
\end{equation}
For any holomorphic scalar $\kappa(t)$, $H_z$ may be replaced in
\eqref{f:A-contour} by $H_z-\kappa H$.  If on the outer wall
$|H|\ge\delta>0$ and $|H_z-\kappa H|\le L_\kappa$, then
\begin{equation}\label{f:A-bound}
          |A|\le K:=\frac{4RL_\kappa}{\delta(R-r)^2}.
\end{equation}
Every zero of $\Delta$ is simple and has derivative $8$.
\end{proposition}

\begin{proof}
At a simple root the heat equation gives $z_i'=H_{zz}/H_z$.
Differentiating $UP$ at that root yields
\[
 z_1'=\frac2{z_1-z_2}+2q(t,z_1),\qquad
 z_2'=\frac2{z_2-z_1}+2q(t,z_2).
\]
Thus
\[
 \Delta'=8+4\Delta\frac{q(t,z_1)-q(t,z_2)}{z_1-z_2}.
\]
By residues the divided difference is
$(2\pi i)^{-1}\oint q/P\,d\zeta$.  Since $f=q+P_z/P$ and
$P_z/P^2=-\partial_z(1/P)$, the integral is unchanged if $q$ is
replaced by $f$.  This proves the equation away from events.
The discriminant is not identically zero by
Proposition~\ref{f:hermite}; alternatively, a persistent double
$H=U(z-c(t))^2$ would give $H_t(t,c(t))=0$ and
$H_{zz}(t,c(t))=2U(t,c(t))\ne0$, contradicting heat flow.
Both sides of \eqref{f:affine-equation} are holomorphic, so the equation
extends through its isolated events.

The scalar subtraction is exact because $\oint1/P\,d\zeta=0$:
the rational function $1/P$ has no exterior finite poles and decays
quadratically at infinity.  On the outer wall,
$|P|\ge(R-r)^2$.  Estimating the contour length $2\pi R$ gives
\eqref{f:A-bound}.  Substitution of $\Delta=0$ in the equation gives
$\Delta'=8$.
\end{proof}

There is no inverse internal root separation in
\eqref{f:A-contour} or \eqref{f:A-bound}.  The bound does pay for
the exterior collar and the spatial-wall floor.  A scalar centering
$\kappa$ improves it only if the corresponding coupled numerator is
actually estimated more sharply.

\begin{theorem}[Reduction of complex charge to two real cap signs]
\label{f:real-cap}
Let
\[
 D=\{x+iy:t_-<x<t_+,\ |y|<\tau\},\qquad t_-<t_+,
\]
and suppose the proper two-root and collar hypotheses hold on a
neighborhood of its closure.  Let $a\in\mathbb R$ and assume reality
$H(\overline t,\overline z)=\overline{H(t,z)}$.  It suffices that
\begin{equation}\label{f:phase}
 \sup_{x+iy\in\overline D}
   \left|\int_0^y\operatorname{Re} A(x+is)\,ds\right|<\frac\pi2.
\end{equation}
In particular, $|A|\le K$ and $K\tau<\pi/2$ suffice.  Then there
is at most one cluster event in $\overline D$, and any such event has
real time and real spatial position.  If
$\Delta(t_-)\Delta(t_+)\ne0$, the entire product boundary is free
of joint zeros, and
\begin{equation}\label{f:cap-charge}
 Q_{D\times\Omega_R}
  =\frac{\operatorname{sgn}\Delta(t_+)
               -\operatorname{sgn}\Delta(t_-)}2\in\{0,1\}.
\end{equation}
\end{theorem}

\begin{proof}
Choose $t_*\in(t_-,t_+)$ and define the holomorphic primitives
\[
 J(t)=\int_{t_*}^t A(v)\,dv,\qquad
 \mu(t)=e^{-J(t)},\qquad \Psi(t)=\int_{t_*}^t\mu(v)\,dv.
\]
They are single valued on the rectangle.  Reality and uniqueness of
the prepared polynomial imply reality of $A$ on the real interval.
Integrating first horizontally and then vertically gives
\[
 \operatorname{Im}J(x+iy)=\int_0^y\operatorname{Re}A(x+is)\,ds.
\]
Therefore \eqref{f:phase} implies $\operatorname{Re}\mu>0$ on the
closed rectangle.  Convexity gives, for distinct $s,t$ in its closure,
\[
 \frac{\Psi(t)-\Psi(s)}{t-s}
      =\int_0^1\mu(s+v(t-s))\,dv,
\]
whose real part is positive.  Thus $\Psi$ is injective, including
on the boundary.  Multiplying \eqref{f:affine-equation} by $\mu$
and integrating yields
\begin{equation}\label{f:normal-form}
                 \mu(t)\Delta(t)=\Delta(t_*)+8\Psi(t).
\end{equation}
There is at most one zero.  A nonreal zero would have a distinct
conjugate zero, so any zero is real.  At a real event the real monic
quadratic has a repeated root, which is necessarily real as well.

The only possible event times on the boundary are now $t_-$ and
$t_+$, and the cap assumptions exclude them.  On real time $\mu>0$
and $(\mu\Delta)'=8\mu>0$.  Equal nonzero cap signs give no event;
negative then positive signs give exactly one.  The reversed sign
pattern cannot occur.  Each event contributes one by
Proposition~\ref{f:affine} and Theorem~\ref{f:charge-winding}, proving
\eqref{f:cap-charge}.
\end{proof}

The theorem proves an upper bound of one event, not zero events.  The
heat family $H=z^2-2t$ has $A=0$ and $\Delta=8t$; on
$|\operatorname{Re}t|,|\operatorname{Im}t|\le1/4$ its roots lie in
$|z|<1$, so $r=1,R=2$ supply a proper collar and charge one.  The
phase condition is entirely compatible with this collision.  This
distinction is central to the endpoint problem \cite{note05}.

\subsection{A complete-source strip criterion}

For the complete positive source, use the moments
\[
 M_j(T,s)=\int_0^\infty\Phi(u)u^j e^{Tu^2}\cosh(su)\,du.
\]
They are finite, and no source modes or exterior tail are omitted.
Suppose $d_r,d_R>0$ are established bounds
\[
 |H(x,z)|\ge d_s\quad
 (t_-\le x\le t_+,\ |z-a|=s),\qquad s=r,R,
\]
and a real anchor has root count two on both circles.  Put
$B_s=M_2(t_+,s)$ and $L=M_1(t_+,R)$.  Then a sufficient width is
\begin{equation}\label{f:source-strip}
 0<\tau<\min\left\{
    \frac{d_r}{B_r},
    \frac{d_R}{B_R+8RL/(\pi(R-r)^2)}\right\}.
\end{equation}
Indeed, $|e^{iyu^2}-1|\le|y|u^2$ and
$|\cos(zu)|,|\sin(zu)|\le\cosh(su)$ on a circle centered at a
real $a$ give
\[
 |H(x+iy,z)-H(x,z)|\le\tau B_s,\qquad
 |H_z(x+iy,z)|\le M_1(t_+,s).
\]
Both walls therefore have positive floors $d_s-\tau B_s$, and
continuity of their counts gives the proper two-root collar throughout
the strip.  Applying \eqref{f:A-bound} with $\kappa=0$ and
$\delta=d_R-\tau B_R$ gives $K\tau<\pi/2$ exactly from the second
inequality in \eqref{f:source-strip}.  Purely imaginary time lifting
requires no extra factor $e^{\tau u^2}$.  The cap signs still require
separate information.  The finite verified application discussed later
uses this type of complete-source payment \cite{note06}.

\section{Explicit windows at a hypothetical double}
\label{f:window-section}

The preceding theorem has a pointwise geometric contract at every
finite real height.  This section pays it in terms of the local
second jet, including when the two roots coincide \cite{note07}.
Fix $T_0,R_0>0$ and set
\begin{equation}\label{f:B-moments}
 B_j=M_j(T_0,R_0),\qquad j=2,3.
\end{equation}
The complete source gives
$|\partial_z^jH(t,z)|\le B_j$ whenever
$\operatorname{Re}t\le T_0$ and $|\operatorname{Im}z|\le R_0$.
These bounds do not depend on $\operatorname{Re}z$.

\begin{theorem}[Quantitative double window]\label{f:double-window}
Assume, conditionally, that for a real $a$,
\[
 H(0,a)=H_z(0,a)=0,\qquad
 c=\tfrac12H_{zz}(0,a)\ne0,
 \qquad q_0=|c|.
\]
Define
\begin{equation}\label{f:window-constants}
 \begin{split}
 R&=\min\{R_0/2,q_0/B_3\},\qquad r=R/2,\\
 h&=\min\left\{T_0/2,
       \frac{q_0R^2}{16\sqrt2 B_2},\frac{R^2}{128}\right\}.
 \end{split}
\end{equation}
For every $t$ with $|\operatorname{Re}t|,|\operatorname{Im}t|\le h$,
$H(t,\cdot)$ has exactly two roots inside each of the spatial circles
of radii $r$ and $R$, and their walls satisfy
\begin{equation}\label{f:double-floors}
 |H(t,z)|>\tfrac12q_0s^2\quad(|z-a|=s,\ s=r,R).
\end{equation}
On the outer wall $|H_z|<3q_0R$.  The regular coefficient satisfies
\begin{equation}\label{f:double-phase}
          |A|\le K:=96/R^2,\qquad Kh\le3/4<\pi/2.
\end{equation}
The hypotheses extend to a neighborhood of the closed product.
\end{theorem}

\begin{proof}
For $|w|=s\le R$, spatial Taylor expansion with integral remainder
and straight-line integration in complex time give
\[
 |H(0,a+w)-cw^2|\le B_3s^3/6\le q_0s^2/6,
\]
\[
 |H(t,a+w)-H(0,a+w)|\le|t|B_2
       \le\sqrt2hB_2\le q_0R^2/16=q_0r^2/4.
\]
At $s=r,R$ their sum is at most $5q_0s^2/12<q_0s^2$.
Rouch\'e's theorem against $cw^2$ gives both counts and the stronger
common floor $7q_0s^2/12$.  At the outer circle the same calculation
gives $37q_0R^2/48$.  Counting twice excludes annular roots.

For the first derivative use its spatial remainder and
$H_{zt}=-H_{zzz}$:
\[
 \begin{split}
 |H_z(t,a+w)|
 &\le2q_0R+\tfrac12 B_3R^2+\sqrt2hB_3\\
 &\le\left(\tfrac52+\frac{\sqrt2}{128}\right)q_0R<3q_0R.
 \end{split}
\]
The last step uses $RB_3\le q_0$ and $h\le R^2/128$.
Equation~\eqref{f:A-bound}, with
$\delta=q_0R^2/2$ and $R-r=R/2$, now gives $K=96/R^2$.
All required wall margins are strict and $Kh\le3/4<\pi/2$;
continuity on the compact boundary preserves the conditions on a
small enlargement.  The corner factor $\sqrt2$ was included explicitly
in the time estimates.
\end{proof}

\begin{corollary}[The compatible cap signs]\label{f:double-caps}
In Theorem~\ref{f:double-window} define, for $x\ge0$,
\begin{equation}\label{f:LU}
 L_K(x)=\frac8K(1-e^{-Kx}),\qquad
 U_K(x)=\frac8K(e^{Kx}-1),
\end{equation}
with continuous limiting values $8x$ if $K=0$.  Then
\begin{equation}\label{f:double-cap-bounds}
 \begin{split}
 L_K(h)&\le\Delta(h)\le U_K(h),\\
 -U_K(h)&\le\Delta(-h)\le-L_K(h).
 \end{split}
\end{equation}
The resulting product has charge one, with its unique event at $(0,a)$.
\end{corollary}

\begin{proof}
At the assumed double, $\Delta(0)=0$.  On real time the affine equation
has the exact solution
\[
 \Delta(t)=8\int_0^t
       \exp\left(\int_s^t A(v)\,dv\right)\,ds.
\]
The coefficient is real and bounded by $K$, so comparison with
$e^{-K|t-s|}$ and $e^{K|t-s|}$ gives
\eqref{f:double-cap-bounds}.  In particular the cap margins are
positive, and $L_K(h)\ge8he^{-Kh}$.  Theorem~\ref{f:real-cap}
gives the charge and uniqueness.
\end{proof}

This is a geometry theorem conditional on an exact double.  It cannot
also be an exclusion theorem: its actual cap signs are those of a
collision.  Improving an absolute coefficient bound changes their
quantitative margins, but cannot turn them into equal signs.

\subsection{A finite-jet test away from an assumed event}

At any real center let
$c_0=H(0,a)$, $c_1=H_z(0,a)$, $c_2=H_{zz}(0,a)/2$, and assume
$q_0=|c_2|>0$.  Choose $0<r<R\le R_0$ and
$\sqrt2h\le T_0$.  If, for $s=r,R$,
\begin{equation}\label{f:off-event}
 d_s=q_0s^2-|c_0|-s|c_1|-B_3s^3/6-\sqrt2hB_2>0,
\end{equation}
then the same Rouch\'e proof gives two roots inside each circle and
wall floors $d_s$ on the closed time square.  Moreover
\begin{equation}\label{f:off-event-A}
 L=|c_1|+2q_0R+B_3R^2/2+\sqrt2hB_3,
 \qquad K=\frac{4RL}{d_R(R-r)^2}
\end{equation}
are valid outer derivative and coefficient bounds.  If $Kh<\pi/2$,
the real-cap theorem applies once the two cap discriminants are
enclosed away from zero.  The test does not assert an event, an
internal gap, or either cap sign.  An interval implementation needs a
lower enclosure for $q_0$ and complete upper enclosures for every
error term; uncertain jet equalities cannot be promoted to an exact
double.

\subsection{Event duration and the height quantifier}

If an event is already known to occur at real time $s<t_*$ and
$|A|\le K$ between them, put $d=t_*-s$ and
$\delta=\Delta(t_*)>0$.  The preceding integral solution gives
$L_K(d)\le\delta\le U_K(d)$.  Therefore, if
$0<\delta_-\le\delta\le\delta_+$ and $K\delta_+<8$,
\begin{equation}\label{f:duration}
 \frac{\log(1+K\delta_-/8)}K
 \le d\le
 -\frac{\log(1-K\delta_+/8)}K.
\end{equation}
For $K=0$ the interval is $[\delta_-/8,\delta_+/8]$.
These inequalities require event existence, supplied for example by
opposite cap signs; a single positive cap value does not supply it.
The scalar bounds are sharp under the information $|A|\le K$, as
constant coefficients $A=K$ and $A=-K$ attain the two extremes.
This statement does not assert that either coefficient is realized by
theta.  A strictly negative event location would follow only from
data making the resulting upper bound for $s$ negative.

The moments in \eqref{f:B-moments} work for every finite real $a$.
Thus every hypothetical exact double has a positive window from
\eqref{f:window-constants}, point by point.  No uniform lower bound
for its second jet is proved.  In fact, with
$f(u)=u^2\Phi(u)$ and $C=\|f'\|_{L^1(0,\infty)}/2<\infty$,
integration by parts gives, for $a\ne0$,
\[
 q_0(a)=\frac12\left|\int_0^\infty f(u)\cos(au)\,du\right|
         \le\frac C{|a|}.
\]
Consequently this particular prescription obeys
\begin{equation}\label{f:height-scale}
 R(a)\le\frac{C}{B_3|a|},\qquad
 h(a)\le\frac{C^3}{16\sqrt2 B_2B_3^2|a|^3}.
\end{equation}
These are upper bounds for conservatively chosen window sizes, not
necessary scales for every possible construction.  On a bounded
real interval there are only finitely many zeros of the nonzero entire
function $H(0,\cdot)$, so the windows of its exact doubles there have
positive minima if the set is nonempty.  This does not compute those
minima or identify any doubles.  An all-height exclusion may work
pointwise and need not use a common negative cap, but its signed
estimate must cover every admissible finite $a$ and the associated
shrinking parameters.  The geometric quantifier is paid here; that
signed quantifier remains open.

\section{Signed source response and displacement of a double event}
\label{s:response}

The two-root flow separates the analytic geometry of a possible collision
from the source information needed to exclude it.  This section gives a
signed source functional, an explicit nonlinear error, and a continuation
criterion for the event itself.  The results consolidate the analytic
arguments in \cite{note08}; their proofs are included here.

\subsection{An exact response formula at a fixed real time}

Fix a real time, which is suppressed in the notation, and a real spatial
center $a$.  Let $F$ and $H=F+E$ be holomorphic near
$\overline{B(a,R)}$.  Assume that $F$ is nonzero on
$\Gamma=\{z:|z-a|=R\}$, has exactly two zeros counted with multiplicity
inside $\Gamma$, and that both zeros belong to $B(a,r)$, where $0<r<R$.
Write $c$ for their centroid and $\Delta_F$ for their squared difference.
The reference zeros may coincide.  Throughout this section contours have
counterclockwise orientation, and
\[
 \mathcal I[g]=\frac{1}{2\pi i}\int_\Gamma g(z)\,dz.
\]
Suppose the complete wall estimate is
\begin{equation}
 \rho:=\sup_\Gamma|E/F|<1.
 \label{s:ratio}
\end{equation}
Rouch\'e's theorem then gives two zeros of $F+sE$ inside $\Gamma$ for
$0\le s\le1$.  A count in a smaller disk, if needed for the two-root
flow theorem, is an additional condition; it does not follow from
\eqref{s:ratio} alone.

\begin{theorem}[Logarithmic discriminant response]
\label{s:log-response}
Let $\ell=\log(1+E/F)$ be the branch defined by its convergent power
series on an annulus containing $\Gamma$.  Then
\begin{equation}
 \Delta_H-\Delta_F
 =-4\mathcal I[(z-c)\ell]-\bigl(\mathcal I[\ell]\bigr)^2.
 \label{s:exact-response}
\end{equation}
Set
\begin{align}
 \mathcal L_F(E)&=-4\mathcal I[(z-c)E/F],\label{s:linear-response}\\
 l_1(\rho)&=-\log(1-\rho),\qquad
 l_2(\rho)=-\log(1-\rho)-\rho,\notag\\
 \mathcal B(\rho)&=4R(R+r)l_2(\rho)+R^2l_1(\rho)^2.
 \label{s:quadratic-payment}
\end{align}
The finite response satisfies
\begin{equation}
 |\Delta_H-\Delta_F-\mathcal L_F(E)|\le\mathcal B(\rho).
 \label{s:response-bound}
\end{equation}
In particular the error is $O(\rho^2)$, with constants independent of
the distance between the two reference zeros.
\end{theorem}

\begin{proof}
Put $w=z-a$, and for $G=F,H$ let
$S_k(G)=\mathcal I[w^kG_z/G]$ be the sum of the $k$th powers of its
two centered zeros, counting multiplicity.  On the annulus,
$H_z/H-F_z/F=\ell_z$.  Integration by parts on the closed contour gives
\[
 S_1(H)-S_1(F)=-\mathcal I[\ell],\qquad
 S_2(H)-S_2(F)=-2\mathcal I[w\ell].
\]
The identity $\Delta_G=2S_2(G)-S_1(G)^2$, together with
$S_1(F)=2(c-a)$, proves \eqref{s:exact-response}.  No extension of
$\ell$ across the zeros of $F$ was used.  Its power series yields
$|\ell|\le l_1(\rho)$ and $|\ell-E/F|\le l_2(\rho)$ on $\Gamma$.
Now $|\mathcal I[g]|\le R\sup_\Gamma|g|$ and
$|z-c|\le R+r$ prove \eqref{s:response-bound}.
\end{proof}

If $F$ and $H$ are real-symmetric, their discriminants and
$\mathcal L_F(E)$ are real.  Thus a signed enclosure
$\mathcal L_F(E)\in[L_-,L_+]$ gives
\begin{equation}
 \Delta_F+L_- -\mathcal B(\rho)
 \le\Delta_H\le
 \Delta_F+L_+ +\mathcal B(\rho).
 \label{s:signed-enclosure}
\end{equation}
A positive lower bound at a negative real cap can be combined with the
two-root charge theorem once all its spatial and phase hypotheses have
been established.  A positive lower bound at $t=0$ itself excludes an
endpoint double in this fiber.  Formula \eqref{s:signed-enclosure}
does not provide either strict inequality without further source
information.  It also retains the external wall ratio and the radii;
being finite at an internal double is not a uniform condition estimate
over all heights.

\subsection{The complete signed source functional}

Suppose, now displaying the fixed real time $t$, that
\[
 E(t,z)=\int_0^\infty\psi(u)e^{tu^2}\cos(zu)\,du.
\]
If $|F(t,z)|\ge\delta>0$ on $\Gamma$ and
\begin{equation}
 N_\psi(t,R):=\int_0^\infty|\psi(u)|e^{tu^2}\cosh(Ru)\,du<\infty,
 \label{s:complete-residual-norm}
\end{equation}
then Fubini's theorem gives
\begin{align}
 \mathcal L_F(E)&=\int_0^\infty\psi(u)e^{tu^2}K_F(u)\,du,
 \label{s:source-functional}\\
 K_F(u)&=-4\mathcal I\!\left[
             \frac{(z-c)\cos(zu)}{F(t,z)}\right],
 \qquad
 |K_F(u)|\le\frac{4R(R+r)}{\delta}\cosh(Ru).
 \label{s:source-kernel}
\end{align}
Indeed $|\operatorname{Im}z|\le R$ on the circle because $a$ is real,
so $|\cos(zu)|\le\cosh(Ru)$.  This proves an integrable majorant for
the exchanged integrals, and also
$\rho\le N_\psi(t,R)/\delta$.  Conjugation symmetry makes $K_F$ real
on the positive real $u$ axis.

At an exact double $b=c$, factor $F(t,z)=U(t,z)(z-b)^2$ with $U$ a
nonvanishing analytic function.  A residue in
\eqref{s:linear-response} then gives the particularly simple identity
\begin{equation}
 \mathcal L_F(E)=-\frac{4E(t,b)}{U(t,b)}
               =-\frac{8E(t,b)}{F_{zz}(t,b)}.
 \label{s:double-response}
\end{equation}
The sign is therefore an oscillatory source measurement divided by a
nonzero curvature.  Positivity of the source and absolute moment bounds
alone do not determine it.  For numerical use, every omitted source
mode, unbounded integration ray, contour arc, and quadrature error must
be included in both the wall ratio and the signed functional.  No new
numerical acquisition is used to prove this identity.

\subsection{Differentiating a real event through an affine heat family}

Let $H_s(t,z)=F(t,z)+sE(t,z)$ be real-symmetric, with both $F$ and $E$
satisfying the backward heat equation.  Suppose $F$ has an ordinary real
double at $(\tau_0,b_0)$, so $F=F_z=0$ and $F_{zz}\ne0$ there.

\begin{proposition}[Event derivatives and a one-sided finite error]
\label{s:event-derivatives}
There is a unique local real-analytic event branch
$s\mapsto(\tau_s,b_s)$.  At a point on this branch, set
\[
 N=H_{s,zz}(\tau_s,b_s),\quad M=H_{s,zzz}(\tau_s,b_s),\quad
 P_4=H_{s,zzzz}(\tau_s,b_s).
\]
With the residual derivatives evaluated at the same point,
\begin{align}
 \tau_s'&=\frac{E}{N}, &
 b_s'&=\frac{M\tau_s'-E_z}{N},\label{s:event-first}\\
 \tau_s''&=\frac{P_4}{N}(\tau_s')^2
          -\frac{2E_{zz}}{N}\tau_s'-(b_s')^2.
 \label{s:event-second}
\end{align}
If the branch exists for $0\le s\le1$ and along it
\[
 |N|\ge\nu>0,\quad |P_4|\le C_4,\quad |E|\le E_0,
 \quad |E_{zz}|\le E_2,
\]
then
\begin{equation}
 \tau_1\le\tau_0+\frac{E(\tau_0,b_0)}{F_{zz}(\tau_0,b_0)}
       +\frac12\left(\frac{C_4E_0^2}{\nu^3}
                         +\frac{2E_2E_0}{\nu^2}\right).
 \label{s:event-location}
\end{equation}
\end{proposition}

\begin{proof}
The Jacobian of $(H_s,H_{s,z})$ with respect to $(t,z)$ at an event is
\[
 \begin{pmatrix}-N&0\\-M&N\end{pmatrix},
\]
whose determinant is $-N^2$.  The analytic implicit-function theorem
gives the local branch, and real symmetry gives its reality.  One
derivative of the two event equations gives
$-N\tau_s'+E=0$ and $-M\tau_s'+Nb_s'+E_z=0$.
For a second derivative of the first equation, use
$H_{s,tt}=P_4$, $H_{s,tz}=-M$, and $E_t=-E_{zz}$ to obtain
\[
 0=-N\tau_s''+P_4(\tau_s')^2-2M\tau_s'b_s'
       +N(b_s')^2-2E_{zz}\tau_s'+2E_zb_s'.
\]
Substitution of $E_z=M\tau_s'-Nb_s'$ leaves
\eqref{s:event-second}.  All terms are real.  In particular
$-(b_s')^2\le0$, whereas $|\tau_s'|\le E_0/\nu$, so
\[
 \tau_s''\le\frac{C_4E_0^2}{\nu^3}
                         +\frac{2E_2E_0}{\nu^2}.
\]
Integrating twice on $[0,1]$ proves \eqref{s:event-location}.
\end{proof}

This calculation also checks the signs in the discriminant formula.
At a double, $\partial_s\Delta_s=-8E/N$ by
\eqref{s:double-response}, while $\partial_t\Delta_s=8$ by the
two-root heat equation.  Differentiating
$\Delta_s(\tau_s)=0$ recovers $\tau_s'=E/N$.

\subsection{Paying continuation and curvature from common walls}

The local implicit-function theorem by itself does not continue an
event to $s=1$.  The following conditions supply that missing step and
the curvature floor in \eqref{s:event-location}.

\begin{proposition}[A complete conditional continuation criterion]
\label{s:continuation}
Let $\mathcal D$ be a closed time rectangle with real caps $t_-<t_+$
and imaginary half-width $\tau$.  Suppose every $H_s=F+sE$,
$0\le s\le1$, has exactly two zeros inside $B(a,r)$ throughout
$\mathcal D$, none in the closed collar to $\Gamma$, a common outer
wall floor $|H_s|\ge\delta>0$, and a common strict phase condition
from the two-root real-cap theorem.  Assume
\[
 \Delta_F(t_-)<0<\Delta_F(t_+),\qquad
 m_{\rm cap}=\min\{-\Delta_F(t_-),\Delta_F(t_+)\}>0.
\]
If the two outer cap circles satisfy $|E/F|\le\rho<1$ and
\begin{equation}
 C_{\rm cap}:=4R(R+r)l_1(\rho)+R^2l_1(\rho)^2<m_{\rm cap},
 \label{s:cap-persistence}
\end{equation}
then a unique real event exists for every $s\in[0,1]$, the events form
one real-analytic branch, and
\begin{equation}
 |H_{s,zz}(\tau_s,b_s)|\ge\frac{2\delta}{(R+r)^2}.
 \label{s:curvature-payment}
\end{equation}
\end{proposition}

\begin{proof}
Apply \eqref{s:exact-response} to $sE$ at each cap.  The bound
$|\ell|\le l_1(\rho)$ gives
$|\Delta_s(t_\pm)-\Delta_F(t_\pm)|\le C_{\rm cap}$.
Thus both cap signs persist.  The two-root real-cap theorem gives
exactly one event for each $s$, located on the real time interval.  It
has spatial multiplicity two and $\partial_t\Delta_s=8$.  The local
branches of Proposition~\ref{s:event-derivatives} agree wherever their
domains overlap by uniqueness.  They extend along the full compact
parameter interval: any limit of event points remains away from the
time caps by the strict cap margins and from the spatial boundary by
the common walls, and is again a double.  The implicit-function theorem
then extends the branch at that limit.

For the curvature, write $H_s=U_sP_s$ in the disk.  On $\Gamma$,
$|P_s|\le(R+r)^2$, hence $|U_s|\ge\delta/(R+r)^2$.  Since $U_s$
has no zeros, the maximum modulus principle applied to $1/U_s$
propagates this bound throughout the disk.  At an event,
$P_s=(z-b_s)^2$ and $H_{s,zz}=2U_s$, proving
\eqref{s:curvature-payment}.
\end{proof}

Here are sufficient reference data for all the common hypotheses.
Suppose $F$ has two zeros inside each of the disks of radii $r,R$
throughout $\mathcal D$, with wall floors $\delta_{F,r},\delta_{F,R}$,
and
\[
 Q_v=\sup_{t\in\mathcal D,\ |z-a|=v}|E/F|<1\quad(v=r,R).
\]
Rouch\'e's theorem supplies both counts for every $H_s$, with floors
$\delta_v=(1-Q_v)\delta_{F,v}$.  Take $\delta=\delta_R$.  If a
complete outer-wall derivative bound is $|F_z|+|E_z|\le L$, the
sufficient common phase inequality is
\begin{equation}
 \frac{4RL}{\delta(R-r)^2}\,\tau<\frac\pi2.
 \label{s:common-phase}
\end{equation}
A more precise common estimate for the regular coefficient $A_s$ may
replace this one.  If the source densities are $\Phi_0$ and $\psi$,
the remaining bounds along the real event branch can be taken to be
\begin{align}
 E_j&=\int_0^\infty|\psi(u)|u^je^{t_+u^2}\,du\quad(j=0,2),\notag\\
 C_4&=\int_0^\infty\bigl(|\Phi_0(u)|+|\psi(u)|\bigr)
                           u^4e^{t_+u^2}\,du.
 \label{s:moment-payments}
\end{align}
The real event positions allow $|\cos(b_su)|\le1$ in these bounds.
All source modes and integration rays are included.

Consequently \eqref{s:event-location}, together with
\eqref{s:cap-persistence}--\eqref{s:moment-payments}, is a complete
conditional criterion for locating the unique event below zero.  It
does not supply a negative right-hand side for the unchanged theta
source.  Both the signed initial velocity and the external wall
margins remain substantive inputs.

\begin{example}[An exact nonlinear and continuation check]
\label{s:quadratic-control}
For
\[
 H_s(t,z)=z^2-2t+s(bz+c),
\]
one has
\[
 \Delta_s(t)=8t+s^2b^2-4sc,\qquad
 \tau_s=\frac{sc}{2}-\frac{s^2b^2}{8},\qquad b_s=-\frac{sb}{2}.
\]
The initial signed response is $-4c$, while the finite correction is
$s^2b^2$.  The logarithmic integral alone in
\eqref{s:exact-response} contributes $2s^2b^2$; its last square is
essential.  The event equations give $\tau_0'=c/2$ and
$\tau_s''=-b^2/4$, as required.

The continuation inequalities can also be checked with strict rational
margins.  Set $r=1$, $R=2$, $|\operatorname{Re}t|,
|\operatorname{Im}t|\le1/4$, $b=1/16$, and $c=\pm1/32$.
Since $2|t|\le\sqrt2/2<3/4$, the inner and outer reference wall
floors exceed $1/4$ and $13/4$.  Thus the relative errors are less
than $3/8$ and $5/104$, respectively, and both counts persist.  The
outer floor of $H_s$ exceeds $99/32$, so
\eqref{s:curvature-payment} permits $\nu=11/16$.
Here $A_s=0$ exactly, and the reference cap discriminants are $-2,+2$.
Since $l_1(5/104)\le5/99$,
\[
 C_{\rm cap}\le\frac{120}{99}+\frac{100}{9801}<2.
\]
For $c=-1/32$, one has $P_4=E_{zz}=0$, so
\eqref{s:event-location} proves $\tau_1\le-1/64<0$; its exact value
is $-1/64-1/2048$.  For $c=+1/32$ the event is positive.  The same
geometric conditions accommodate both outcomes; the sign comes from
the signed displacement.
\end{example}

\section{What finite complete envelope moments cannot determine}
\label{s:moment-obstruction}

The obstruction proved here concerns precisely the envelope moments
used in absolute source estimates.  It does not apply to an arbitrary
finite list of observations, which could include the oscillatory
measurement one wishes to determine.

\begin{theorem}[Fixed envelope moments with both event directions]
\label{s:finite-moment-control}
Fix any finite collection
\begin{equation}
 k_j(u)=u^{n_j}e^{T_ju^2}\cosh(Y_ju),\qquad
 n_j\in\mathbb Z_{\ge0},\quad T_j\in\mathbb R,\quad Y_j\ge0.
 \label{s:envelope-moments}
\end{equation}
There is a strictly positive smooth even source $\Phi_0$, with
super-Gaussian decay, and a real smooth perturbation $\psi$ supported
in a compact subinterval of $(0,\infty)$, such that for all sufficiently
small positive and negative $s$:
\begin{enumerate}
 \item $\Phi_s=\Phi_0+s\psi$ remains positive, with the same source
       near zero and outside a fixed compact interval;
 \item every selected complete moment
       $\int_0^\infty\Phi_s(u)k_j(u)\,du$ is independent of $s$;
 \item the associated heat family has a real ordinary-double branch
       with $\tau_0=0$ and $\tau_0'>0$.
\end{enumerate}
In particular the double event lies on either side of zero for the
respective signs of sufficiently small $s$.
\end{theorem}

\begin{proof}
Take $c_0>0$ and
\[
 \Phi_0(u)=c_0\left(e^{-u^4}
          +\tfrac12e^{-(u-4)^4}+\tfrac12e^{-(u+4)^4}\right).
\]
This is positive, even, and rapidly decaying.  If
$G(z)=\int_0^\infty e^{-u^4}\cos(zu)\,du$, translation of the
whole-line Fourier integral gives
\[
 H_0(0,z)=c_0G(z)(1+\cos4z).
\]
At $b=\pi/4$ the last factor has an ordinary double.  To see that
$G(b)>0$ without numerical evaluation, write
$M_j=\int_0^\infty u^je^{-u^4}\,du$.  Integration of
$(ue^{-u^4})'$ gives $M_4=M_0/4$; Cauchy--Schwarz gives
$M_2\le(M_0M_4)^{1/2}=M_0/2$.  Since
$\cos x\ge1-x^2/2$,
\[
 G(b)\ge M_0-\frac{b^2}{2}M_2
       \ge M_0(1-b^2/4)>0.
\]
Consequently $H_{0,zz}(0,b)=16c_0G(b)>0$.

Choose a bounded open interval $I\Subset(0,\infty)$ and a smooth
nonnegative bump $\chi$ supported in $I$ and positive on a subinterval.
In the real inner product
$\langle f,g\rangle_\chi=\int_0^\infty\chi fg$, let $p$ be the
orthogonal projection of $f(u)=\cos(bu)$ onto the span of the $k_j$;
linearly dependent functions may first be discarded.  Set
$\psi=\chi(f-p)$.  Then
\begin{equation}
 \int_0^\infty\psi k_j=0,
 \qquad
 \int_0^\infty\psi(u)\cos(bu)\,du
    =\int_0^\infty\chi(u)(f(u)-p(u))^2\,du>0.
 \label{s:projection-control}
\end{equation}

For the strict inequality, $\cos(bu)$ cannot belong to the indicated
finite span on any interval.  Such an equality would extend to the
whole real axis by analyticity.  Expanding each hyperbolic cosine and
grouping equal exponential factors writes its proposed right-hand side
as a finite sum of polynomial multiples of $e^{Tu^2+Yu}$ with real
$T,Y$.  As $u\to+\infty$, eliminate the largest quadratic exponent
$T>0$, and within it the largest linear exponent, since a nonzero
polynomial times that factor cannot remain bounded.  Repeating removes
all growing terms, including those with $T=0,Y>0$.  All terms with
$T<0$, or with $T=0,Y<0$, decay.  The remaining term is a polynomial,
which must be constant to remain bounded.  The resulting right-hand
side has a limit, whereas $\cos(bu)$ does not.  This contradiction
proves \eqref{s:projection-control}.

The base source has a positive minimum on $\operatorname{supp}\psi$;
therefore both signs of sufficiently small $s$ preserve positivity.
The even extension of $\psi$ is smooth because its support is away
from zero, and all prescribed moments remain exactly fixed.  The
super-Gaussian tail and the compact perturbation give jointly entire
heat transforms.  Proposition~\ref{s:event-derivatives} and
\eqref{s:projection-control}, divided by the positive curvature,
give $\tau_0'>0$.  The local implicit-function theorem supplies the
branch and the asserted signs.
\end{proof}

A sufficiently small proper two-root time-space window exists around
the base double: isolate the zero at time zero, use nonvanishing on
nested spatial circles to preserve the two counts, and shrink the time
rectangle until the regular coefficient satisfies the strict phase
condition.  Its strict margins persist for all sufficiently small
$|s|$ by uniform convergence on compact sets.  Thus the two event
directions coexist with the same selected moments and the same strict
local analytic requirements.

\begin{remark}[Scope of the control]
The preserved moment vector is that of the constructed source
$\Phi_0$, not the theta source.  Compact smooth perturbations need not
preserve analyticity in the source variable, stronger shape conditions,
or additional primitive identities.  This construction asserts no
global all-real heat slice, no all-positive-time real-rootedness, no
canonical theta coefficient expansion, and no ordinary Dirichlet
series or full normalized arithmetic reflection identity.  It excludes
only a universal inference from the displayed finite envelope moments
and strict local analytic data.  A measured signed functional such as
\eqref{s:source-functional}, exact arithmetic data, and further
consequences of RH lie outside that information class.  In particular,
this preservation list must not be combined with the distinct lists
belonging to the other collision controls.
\end{remark}

\section{The canonical theta tail in the discriminant response}
\label{s:canonical-tail}

The actual theta tail has more structure than an arbitrary positive
source bounded between two envelopes.  We construct that structure
explicitly, following the source identity in \cite{tail04}, and then
apply it to the signed kernel.  The resulting interval is strictly
narrower than a specified pointwise-envelope interval, but is not
proved to exclude zero.

\subsection{An explicit positive Laplace law}

Retain the complete mode normalization
\[
 \lambda_n=\pi n^2,\qquad
 \phi_n(u)=\lambda_ne^{5u}(2\lambda_ne^{4u}-3)
                              e^{-\lambda_ne^{4u}}.
\]
Every mode is positive for $u\ge0$.  Fix an integer $m\ge1$, and set
\begin{gather}
 S_m=\sum_{n=1}^m\phi_n,\qquad R_m=\sum_{n>m}\phi_n,\notag\\
 \alpha=\lambda_m,\quad\beta=\lambda_{m+1},\quad d=\beta-\alpha,
 \quad v=e^{4u}-1,\quad D_0=2\alpha-3,\quad
 \gamma=\frac{D_0}{2\alpha}>0.\notag
\end{gather}
For $n>m$, define the nonnegative quantities
\begin{equation}
 d_n=\lambda_n-\alpha,\qquad e_n=\lambda_n-\beta,
 \qquad
 a_n=\frac{\lambda_n^2}{\alpha^2}e^{-d_n},\qquad
 b_n=\frac{3\lambda_nd_n}{\alpha^2D_0}e^{-d_n}.
 \label{s:tail-masses}
\end{equation}

\begin{proposition}[Canonical tail measure]
\label{s:laplace-law}
Let $\mathsf X_\gamma$ have exponential density
$\gamma e^{-\gamma x}$ on $[0,\infty)$, and put
\begin{align}
 E_m&=\sum_{n>m}(a_n+b_n)=\frac{R_m(0)}{\phi_m(0)},\notag\\
 \nu_m&=\frac1{E_m}\sum_{n>m}
       \left[a_n\delta_{e_n}
                +b_n\operatorname{Law}(e_n+\mathsf X_\gamma)\right],
 \label{s:tail-probability}\\
 p_m(v)&=\int_0^\infty e^{-xv}\,d\nu_m(x).
 \label{s:laplace-transform}
\end{align}
Then $\nu_m$ is a probability measure and
\begin{equation}
 R_m(u)=E_m\phi_m(u)e^{-dv}p_m(v).
 \label{s:exact-tail}
\end{equation}
If $A=a_{m+1}$ and $\eta=A/E_m$, then
\[
 0<\eta<1,\qquad p_m(0)=1,\qquad p_m(\infty)=\eta,
 \qquad (-1)^jp_m^{(j)}(v)\ge0\quad(j\ge0,v\ge0).
\]
In particular
\begin{equation}
 A\phi_m(u)e^{-dv}\le R_m(u)\le E_m\phi_m(u)e^{-dv}.
 \label{s:tail-envelope}
\end{equation}
\end{proposition}

\begin{proof}
Division of the source polynomials, retaining their lower-order terms,
gives
\begin{equation}
 \frac{\phi_n(u)}{\phi_m(u)}
   =e^{-d_nv}\left(a_n+b_n\frac{\gamma}{\gamma+v}\right).
 \label{s:mode-division}
\end{equation}
For example, with $q=1+v$ the needed algebra is
\[
 \frac{2\lambda_nq-3}{2\alpha q-3}
 =\frac{\lambda_n}{\alpha}
    +\frac{3(\lambda_n-\alpha)}{\alpha(2\alpha q-3)}.
\]
Also $\int_0^\infty e^{-xv}\gamma e^{-\gamma x}\,dx
=\gamma/(\gamma+v)$.  Factor $e^{-dv}$ out of
\eqref{s:mode-division} and sum over all $n>m$ to obtain
\eqref{s:exact-tail} and \eqref{s:tail-probability}.  Evaluating at
$v=0$ proves the expression for $E_m$ and total mass one.  The masses
have Gaussian decay in $n$, and the continuous pieces have exponential
tails, so every moment of $\nu_m$ is finite.  Differentiation under its
integral proves the alternating derivative inequalities, including at
zero.  Its only mass at zero is $a_{m+1}/E_m$; the exponential pieces
have no atoms, and $b_{m+1}>0$, so this number lies strictly between
zero and one.  Dominated convergence gives $p_m(\infty)=\eta$.
The last enclosure follows.
\end{proof}

For later use, $E_m/A$ is bounded as $m$ varies and tends to one.  Here
is a direct estimate.  Put $k=m+1$ and
\[
 \varrho_m=\left(\frac{k+1}{k}\right)^4e^{-\pi(2k+1)}<1.
\]
Successive ratios of $a_n$ for $n\ge k$ are at most $\varrho_m$,
and $b_n/a_n=3d_n/(\lambda_nD_0)\le3/D_0$.  Consequently
\begin{equation}
 1\le\frac{E_m}{A}
 \le\frac{1+3/D_0}{1-\varrho_m}\longrightarrow1.
 \label{s:tail-mass-ratio}
\end{equation}
This also supplies a complete elementary majorant for the infinite
mass sum.

\subsection{A cumulative-integral enclosure for the actual response}

At a fixed real cap take the \emph{unshifted} reference and residual
\begin{equation}
 F(t,z)=\int_0^\infty S_m(u)e^{tu^2}\cos(zu)\,du,
 \qquad E(t,z)=\int_0^\infty R_m(u)e^{tu^2}\cos(zu)\,du.
 \label{s:unshifted-head}
\end{equation}
Assume $F$ has the proper two-root geometry of Section~\ref{s:response}
and outer wall floor $\delta>0$.  Define
\begin{align}
 W(u)&=e^{tu^2}\phi_m(u)e^{-d(e^{4u}-1)},\notag\\
 z_m(u)&=\frac{p_m(e^{4u}-1)-\eta}{1-\eta},\qquad
 d\sigma_m(w)=-z_m'(w)\,dw.
 \label{s:monotone-tail}
\end{align}
The function $z_m$ decreases from one to zero.  Hence $\sigma_m$ is
a probability measure, and $z_m(u)=\sigma_m([u,\infty))$.
The source identity becomes
\begin{equation}
 e^{tu^2}R_m(u)=W(u)\{A+(E_m-A)z_m(u)\}.
 \label{s:affine-multiplier}
\end{equation}
For the real kernel in \eqref{s:source-kernel}, put
\[
 J_K=\int_0^\infty W(u)K_F(u)\,du,
 \qquad C_K(w)=\int_0^wW(u)K_F(u)\,du,
 \qquad C_K(\infty)=J_K.
\]
The bound in \eqref{s:source-kernel} and the complete source decay
make these integrals absolutely convergent.

\begin{theorem}[Canonical signed discriminant interval]
\label{s:tail-discriminant}
With the preceding notation,
\begin{equation}
 \mathcal L_F(E)=A J_K+(E_m-A)
                        \int_0^\infty C_K(w)\,d\sigma_m(w).
 \label{s:tail-response}
\end{equation}
Suppose a complete wall-ratio payment is
\begin{equation}
 \overline\rho=\frac{E_m}{\delta}
                \int_0^\infty W(u)\cosh(Ru)\,du<1.
 \label{s:tail-ratio}
\end{equation}
Then
\begin{align}
 \Delta_H&\ge\Delta_F+A J_K+(E_m-A)\min_{0\le w\le\infty}C_K(w)
                         -\mathcal B(\overline\rho),\notag\\
 \Delta_H&\le\Delta_F+A J_K+(E_m-A)\max_{0\le w\le\infty}C_K(w)
                         +\mathcal B(\overline\rho).
 \label{s:canonical-interval}
\end{align}
The extrema exist on the compactified half-line.  A sharper established
wall ratio can replace \eqref{s:tail-ratio}.
\end{theorem}

\begin{proof}
Substitute \eqref{s:affine-multiplier} into
\eqref{s:source-functional}, write
$z_m(u)=\int\mathbf1_{u\le w}\,d\sigma_m(w)$, and use Fubini.
Absolute integrability follows from $\sigma_m$ having mass one and
$\int W|K_F|<\infty$.  This proves \eqref{s:tail-response}.
Continuity of $C_K$ and its limit at infinity give a compact range,
so its probability average lies between its minimum and maximum.
The upper envelope in \eqref{s:tail-envelope} proves
$\sup_\Gamma|E/F|\le\overline\rho$.  Finally
\eqref{s:response-bound}, with the increasing function
$\mathcal B$, proves \eqref{s:canonical-interval}.
\end{proof}

This enclosure retains every theta mode and the entire source ray.
It is an analytic estimate, independent of a sampling scheme.  An
implementation must additionally enclose the acquisition errors and
the full-domain extrema of $C_K$; no positive numerical lower endpoint
is asserted here.

\subsection{Strict improvement over a matched pointwise envelope}

The following comparison keeps the same known lower and upper source
bounds, the same wall ratio, and the same nonlinear response error.
The pointwise comparison drops only the monotonicity of $z_m$.

\begin{lemma}[Oscillation of the two-root response kernel]
\label{s:kernel-oscillation}
If both reference roots lie in $\operatorname{Re}z>0$, then $K_F$
has strictly positive and strictly negative intervals arbitrarily far
along the positive $u$ axis.
\end{lemma}

\begin{proof}
Factor $F=UP$.  For distinct roots $z_1,z_2$, their centroid satisfies
$(z_i-c)/P_z(z_i)=1/2$.  The residue theorem therefore gives
\begin{equation}
 K_F(u)=-2\left(\frac{\cos(z_1u)}{U(z_1)}
                         +\frac{\cos(z_2u)}{U(z_2)}\right).
 \label{s:kernel-residues}
\end{equation}
At a double $b>0$ the formula is
$K_F(u)=-4\cos(bu)/U(b)$, which has the asserted signs.
For distinct real positive roots, \eqref{s:kernel-residues} is a
nonzero real combination of two distinct positive-frequency cosines.
Its primitive is bounded and its mean square is positive: integrating
the cosine products shows that the time average of the square tends
to half the sum of the squared coefficients.  If the kernel were
eventually nonnegative, its integral on the remaining ray would be
bounded, whereas $K_F^2\le\|K_F\|_\infty K_F$ would then bound the
square integral, contradicting its positive linear growth.  Eventual
nonpositivity is excluded in the same way.  Continuity converts the
resulting arbitrarily distant points of either sign into intervals.

The other possibility under real symmetry is a conjugate pair
$b\pm iy$ with $b,y>0$.  Write $1/U(b+iy)=c+ie$.  Equation
\eqref{s:kernel-residues} becomes
\[
 K_F(u)=-4\{c\cos(bu)\cosh(yu)+e\sin(bu)\sinh(yu)\}.
\]
After division by $e^{yu}$ this is a nonzero sinusoid plus a term
tending exponentially to zero; $(c,e)\ne(0,0)$ because $U$ is
nonvanishing.  Its positive and negative lobes give the claim.
\end{proof}

\begin{theorem}[Strict matched-envelope improvement]
\label{s:strict-improvement}
Under the hypotheses of Theorem~\ref{s:tail-discriminant}, assume also
that both reference roots lie in $\operatorname{Re}z>0$.  Set
\[
 f(u)=W(u)K_F(u),\qquad
 P_+=\int_0^\infty f_+(u)\,du,
 \qquad P_-=\int_0^\infty f_-(u)\,du,
\]
where $f_+=\max(f,0)$ and $f_-=\max(-f,0)$.
The pointwise relaxation using the same bounds
$AW\le e^{tu^2}R_m\le E_mW$ has the exact interval for the linear response
\begin{equation}
 [A J_K-(E_m-A)P_-,\ A J_K+(E_m-A)P_+].
 \label{s:pointwise-interval}
\end{equation}
There are explicit numbers $d_-,d_+>0$, given by integrals over
ordered sign intervals, for which
\begin{equation}
 \min C_K\ge-P_-+d_-,\qquad
 \max C_K\le P_+-d_+.
 \label{s:ordered-gains}
\end{equation}
Thus \eqref{s:canonical-interval} strictly improves both endpoints
of the discriminant interval obtained from
\eqref{s:pointwise-interval} with the identical nonlinear payment
$\mathcal B(\overline\rho)$.
\end{theorem}

\begin{proof}
The pointwise relaxation permits an arbitrary measurable multiplier
$0\le a(u)\le1$ in $W\{A+(E_m-A)a\}$.  Choosing $a=1$ exactly
where $f$ is positive or negative proves the two extrema in
\eqref{s:pointwise-interval}; intermediate values follow by convexity.

Since $W>0$, Lemma~\ref{s:kernel-oscillation} supplies a compact
negative interval $I_-$ followed by a compact positive interval $I_+$.
Set
\[
 d_+=\min\left\{\int_{I_-}f_-\,du,
                         \int_{I_+}f_+\,du\right\}>0.
\]
For every $w$, the identity
\[
 P_+-C_K(w)=\int_w^\infty f_+(u)\,du
                              +\int_0^w f_-(u)\,du
\]
is at least $d_+$: before the end of $I_-$ the later positive interval
is wholly omitted, and after that point the negative interval is
wholly included.  This proves the upper inequality.  A positive
interval followed by a negative one gives, by the same argument,
\[
 d_-:=\min\left\{\int_{I_+^{\prime}}f_+\,du,
                         \int_{I_-^{\prime}}f_-\,du\right\}>0,
 \qquad C_K(w)+P_-\ge d_-.
\]
Finally $E_m-A>0$, so both improvements remain strict after adding
the same respective nonlinear errors.
\end{proof}

The right-half-plane condition is material: a reference on the
imaginary axis can have a kernel of one sign.  The strict comparison
just proved can be arbitrarily small and its improved interval may
still contain zero.  It establishes neither a uniform useful gain as
height increases nor optimality among estimators retaining the full
Laplace law.  It compares two specified source relaxations, not all
methods of arithmetic evaluation.

\subsection{Conditional applicability and the remaining source issue}

\begin{proposition}[Local applicability at every putative positive double]
\label{s:tail-local-applicability}
Suppose the unchanged theta field has an ordinary double at $(0,a)$
with $a>0$.  On a sufficiently small fixed proper two-root window about
that event, the unshifted reference \eqref{s:unshifted-head} satisfies
the hypotheses of Theorems~\ref{s:tail-discriminant} and
\ref{s:strict-improvement} for every sufficiently large cutoff $m$.
The assertion is pointwise in the event; it supplies no uniform cutoff
over all heights and no contradiction to the assumed double.
\end{proposition}

\begin{proof}
In the quantitative double-window construction, impose additionally
$R\le a/2$ and then choose the inner radius and time widths with the
same inequalities.  The proof of its positive wall margins is
unchanged.  Every point in the spatial disk has positive real part.
On this fixed compact window the complete theta tails tend uniformly
to zero on both nested walls.  For all sufficiently large $m$, their
absolute bounds are less than half the corresponding wall floors of
$H$.  Rouch\'e's theorem gives two zeros of $F=H-E$ inside each disk,
none in the collar, and an outer reference floor at least half that of
$H$.  In particular the reference roots lie in the right half-plane,
and $|E/F|\le|E|/(|H|-|E|)<1$ on the outer wall.

The explicit sufficient ratio \eqref{s:tail-ratio} also tends to zero.
By \eqref{s:tail-envelope},
\[
 E_m\int_0^\infty W(u)\cosh(Ru)\,du
 \le\frac{E_m}{A}
       \int_0^\infty R_m(u)e^{tu^2}\cosh(Ru)\,du.
\]
The mass ratio is bounded by \eqref{s:tail-mass-ratio}; the complete
tail integral tends to zero uniformly on the compact time interval by
dominated convergence.  The reference floor stays bounded below, so
\eqref{s:tail-ratio} eventually holds.  The cited theorems now apply.
\end{proof}

The reference in \eqref{s:unshifted-head} is essential.  A shifted or
reflection-corrected head changes the residual; the corresponding
connector terms and shifted integration rays must then be restored
before applying a source identity.  The present proof uses the
unshifted head and its complete tail, and makes no claim to resolve
their small signed cancellation.  This distinction is developed in
\cite{tail05}.

There is also a limitation within the weighted-square source class.
Positive weights for which the relevant source sums and integrals
converge leave the same construction with masses $a_n,b_n$ multiplied
by those weights.  Thus positivity of the Laplace
law does not characterize the canonical coefficient sequence.  When
another source agrees through mode $m+1$ and has the same source value
at zero, it has the same $F,W,A,E_m$ and hence the same monotone
response enclosure at any fixed admissible reference window.  Its
actual probability law and response need not agree with theta's.
No assertion that a prescribed window contains a collision of that
other source follows from such a match.  This observation is separate
from the finite envelope-moment control in
Theorem~\ref{s:finite-moment-control}; their preservation hypotheses
are different.

The analytic output is therefore the actual-source enclosure
\eqref{s:canonical-interval}, its strict matched comparison
\eqref{s:ordered-gains}, and its conditional applicability near each
putative ordinary double.  The sign of the full theta discriminant
lower bound remains open.  The proof and control boundaries were
internally checked in \cite{review08}; these statements do not claim
independent specialist validation.

\section{Discriminant evolution for a finite cluster of arbitrary size}
\label{m:section}

The two-root equation has a regular extension to every finite cluster,
but its constant forcing does not.  This distinction gives both a useful
general identity and exact limits on what can be inferred from two real
cap signs.  The results in this section develop the analytic statements
recorded in \cite{note09}; their earlier internal review is \cite{review08}.

\subsection{A regular equation through collisions}

Let $D$ be a bounded Jordan complex-time domain with piecewise smooth
boundary, and let $H$
be holomorphic near
$\overline D\times\{z:|z-a|\leq R\}$, with $H_t=-H_{zz}$.
Assume that, for every $t\in\overline D$, $H(t,\cdot)$ has exactly
$m\geq2$ zeros in the outer disk, all in $|z-a|<r<R$, and has no zero
on the outer circle.  These assumptions are understood with strict
margins sufficient to extend them to a neighborhood of the closed
product.  The cluster factorization gives
\begin{equation}\label{m:factor}
 P(t,z)=\prod_{i=1}^m(z-z_i(t)),\qquad H(t,z)=U(t,z)P(t,z),
 \qquad q=\frac{U_z}{U},\qquad f=\frac{H_z}{H},
\end{equation}
where $P$ has single-valued holomorphic coefficients and $U$ is
holomorphic and nonvanishing on the spatial disk.  Individual root
labels need not be single-valued.  In particular, $q$ includes the
complete exterior contribution, including any zero-free factor of $H$.

Put $N=m(m-1)/2$ and define
\begin{equation}\label{m:DT}
 \mathfrak D=\prod_{i<j}(z_i-z_j)^2,\qquad
 \mathfrak T=\sum_{i<j}
       \prod_{\substack{k<\ell\\(k,\ell)\ne(i,j)}}
                         (z_k-z_\ell)^2.
\end{equation}
The empty product is one.  Both expressions are symmetric polynomials
in the roots and hence are holomorphic functions of the coefficients
of $P$.  Thus \eqref{m:DT} defines $\mathfrak T$ at collisions; the
quotient expression
$\mathfrak T=\mathfrak D\sum_{i<j}(z_i-z_j)^{-2}$ is used only when
the roots are distinct.

\begin{theorem}[Regular finite-cluster discriminant equation]
\label{m:evolution}
Under these assumptions,
\begin{align}
 \mathfrak D'&=8\mathfrak T+B\mathfrak D,
                                                        \label{m:ode}\\
 B(t)&=\frac{2}{2\pi i}\oint_{|\zeta-a|=R}
        \frac{H_z(t,\zeta)}{H(t,\zeta)}
        \frac{P_{zz}(t,\zeta)}{P(t,\zeta)}\,d\zeta.
                                                        \label{m:B}
\end{align}
The coefficient $B$ is holomorphic through every internal
collision.  For any holomorphic scalar $\kappa(t)$, replacing $H_z$
by $H_z-\kappa H$ in \eqref{m:B} leaves it unchanged.  If, on the
complete outer wall,
\begin{equation}\label{m:wallinputs}
 |H|\geq\delta>0,\qquad |H_z-\kappa H|\leq L_\kappa,
\end{equation}
then
\begin{equation}\label{m:Bbound}
 |B|\leq
 \frac{2R\,m(m-1)L_\kappa}{\delta(R-r)^2}.
\end{equation}
In particular, the bound contains no inverse internal root separation.
\end{theorem}

\begin{proof}
At a simple zero, differentiation and the heat equation give
\begin{equation}\label{m:motion}
 z_i'=\frac{H_{zz}(t,z_i)}{H_z(t,z_i)}
      =2\sum_{j\ne i}\frac{1}{z_i-z_j}+2q(t,z_i).
\end{equation}
Writing $d_{ij}=z_i-z_j$, we obtain
\begin{equation}\label{m:pairmotion}
 \frac{z_i'-z_j'}{d_{ij}}
 =\frac{4}{d_{ij}^2}
 -2\sum_{k\ne i,j}\frac{1}{(z_i-z_k)(z_j-z_k)}
 +2\frac{q(t,z_i)-q(t,z_j)}{d_{ij}}.
\end{equation}
The middle terms cancel after summation over all pairs, because every
unordered triple contributes a multiple of the identity
\[
 \frac{1}{(x-z)(y-z)}+
 \frac{1}{(x-y)(z-y)}+
 \frac{1}{(y-x)(z-x)}=0.
\]
Logarithmic differentiation therefore yields
\begin{equation}\label{m:logderivative}
 \frac{\mathfrak D'}{\mathfrak D}
 =8\sum_{i<j}\frac{1}{d_{ij}^2}
  +4\sum_{i<j}\frac{q(t,z_i)-q(t,z_j)}{z_i-z_j}.
\end{equation}
The divided differences in the last term extend holomorphically when
their arguments coincide.  Cauchy's formula and
\[
 \frac{P_{zz}(t,\zeta)}{P(t,\zeta)}
    =2\sum_{i<j}\frac{1}{(\zeta-z_i)(\zeta-z_j)}
\]
give
\[
 4\sum_{i<j}\frac{q(t,z_i)-q(t,z_j)}{z_i-z_j}
 =\frac{2}{2\pi i}\oint_{|\zeta-a|=R}
             q(t,\zeta)\frac{P_{zz}(t,\zeta)}{P(t,\zeta)}\,d\zeta.
\]
Replacing $q$ by $f-P_z/P$ introduces the rational integral of
$P_zP_{zz}/P^2$.  All finite poles of this rational function lie inside
the contour, and it is $O(\zeta^{-3})$ at infinity.  Its total residue
is consequently zero.  This proves \eqref{m:B}.  The same reasoning
applied to $P_{zz}/P=O(\zeta^{-2})$ proves the exact scalar subtraction.

Every denominator in \eqref{m:B} is nonzero on the integration circle,
so the coefficient is holomorphic in time.  The simple-root set is
dense: at any time, the local heat splitting of a zero of multiplicity
$k$ has distinct Hermite leading roots after a nonzero small time
increment.  The finitely many clusters in that fiber can all be
separated by one sufficiently small increment.  Thus $\mathfrak D$
is not identically zero.  Multiplying \eqref{m:logderivative} by
$\mathfrak D$ proves \eqref{m:ode} on a dense set, and then everywhere
by holomorphic continuation.

Finally, $|\zeta-z_i|\geq R-r$ on the circle, whence
\[
 \left|\frac{P_{zz}}{P}\right|
       \leq\frac{m(m-1)}{(R-r)^2}.
\]
Using the contour length $2\pi R$ and \eqref{m:wallinputs} in
\eqref{m:B} proves \eqref{m:Bbound}.
\end{proof}

For $m=2$ one has $\mathfrak T=1$, so the theorem is exactly the
two-root evolution equation.  For larger clusters, it keeps every
root in the birth while retaining the same power two of the spatial
collar in the coefficient bound.

\subsection{An explicit window at a hypothesized full birth}

For the complete theta source, write
\[
 B_j=M_j(T_0,R_0)
   =\int_0^\infty\Phi(u)u^j e^{T_0u^2}\cosh(R_0u)\,du,
 \qquad T_0,R_0>0.
\]
These finite positive moments bound all derivatives needed below.
\begin{proposition}[Parameterized geometry at multiplicity $m$]
\label{m:window}
Suppose $a\in\mathbb R$ and $H(0,\cdot)$ has exact multiplicity
$m\geq2$ at $a$.  Set
\begin{align}
 q_m&=\frac{|\partial_z^mH(0,a)|}{m!}>0,
 &R&=\min\left\{\frac{R_0}{2},
                \frac{(m+1)!q_m}{4B_{m+1}}\right\},
 &r&=\frac R2,                                             \label{m:radii}\\
 h&=\min\left\{\frac{T_0}{2},
                \frac{q_mr^m}{4\sqrt2 B_2}\right\}.
                                                               \label{m:width}
\end{align}
For $|\operatorname{Re}t|,|\operatorname{Im}t|\leq h$, each of the
two disks of radii $r$ and $R$ contains exactly $m$ roots, and
\begin{equation}\label{m:paidwall}
 |H(t,z)|\geq\frac{q_ms^m}{2}
       \quad\hbox{if } |z-a|=s,\quad s\in\{r,R\}.
\end{equation}
The complete coefficient in \eqref{m:B} obeys
\begin{equation}\label{m:paidB}
 |B|\leq
       \frac{16m(m-1)B_1}{q_mR^{m+1}}.
\end{equation}
\end{proposition}

\begin{proof}
Put $c_m=\partial_z^mH(0,a)/m!$.  Spatial Taylor expansion, with the
complete derivative bound on the segment from $a$ to $z$, gives
\[
 |H(0,z)-c_m(z-a)^m|
 \leq\frac{B_{m+1}s^{m+1}}{(m+1)!}
 \leq\frac{q_ms^m}{4} \qquad (|z-a|=s\in\{r,R\}).
\]
Along the straight time segment, $|H_t|=|H_{zz}|\leq B_2$ and
$|t|\leq\sqrt2h$, so
\[
 |H(t,z)-H(0,z)|\leq\sqrt2hB_2\leq\frac{q_mr^m}{4}.
\]
The total error is at most $q_ms^m/2$, strictly less than
$|c_m(z-a)^m|$.  Rouch\'e's theorem gives both root counts and the
triangle inequality gives \eqref{m:paidwall}.  The inequalities have
strict slack relative to the zero-avoidance and root-count requirements,
so the contract extends near the closed product.  In
\eqref{m:Bbound}, take $\kappa=0$, $L_\kappa=B_1$,
$\delta=q_mR^m/2$ and $R-r=R/2$, obtaining \eqref{m:paidB}.
\end{proof}

This construction is conditional on the event and its nonzero order-$m$
derivative.  It gives a positive window at every such finite event,
with all theta modes retained in $B_j$.  It supplies neither a sign of
the discriminant nor nonvanishing on the complex time boundary.
The dependence on $q_m$ and $m$ is explicit and may deteriorate with
height.

\subsection{Birth order, forcing, and charge}

Suppose the entire cluster coalesces at $(t_0,z_0)$.  Write
$s=t-t_0$ and $s=\rho^2$.  The local heat splitting is
\begin{equation}\label{m:hermitesplit}
 z_i(t_0+\rho^2)=z_0+\sqrt2\lambda_i\rho+O(\rho^2),
\end{equation}
where $\lambda_1,\ldots,\lambda_m$ are the distinct real roots of the
monic probabilists' Hermite polynomial $\operatorname{He}_m$.
Indeed, the leading part of parabolic weight $m$ is
\[
 e^{-s\partial_z^2}(z-z_0)^m
    =(2s)^{m/2}\operatorname{He}_m
                   \left(\frac{z-z_0}{\sqrt{2s}}\right).
\]
All higher initial spatial Taylor terms have larger parabolic weight;
analytic continuation of the simple scaled roots gives
\eqref{m:hermitesplit}.

\begin{corollary}[Order at a full birth]\label{m:birth}
At a full multiplicity-$m$ birth,
\begin{align}
 \mathfrak D(t)&=d_ms^N+O(s^{N+1}),
 &d_m&=2^N\prod_{i<j}(\lambda_i-\lambda_j)^2>0,
                                                               \label{m:Dleading}\\
 \mathfrak T(t)&=\frac{Nd_m}{8}s^{N-1}+O(s^N).
                                                               \label{m:Tleading}
\end{align}
If $\mathfrak D$ is nonzero on $\partial D$, the product-domain
charge is
\begin{equation}\label{m:charge}
 Q=\frac{1}{2\pi i}\oint_{\partial D}
                   \frac{\mathfrak D'}{\mathfrak D}\,dt
  =\frac{8}{2\pi i}\oint_{\partial D}
                   \frac{\mathfrak T}{\mathfrak D}\,dt.
\end{equation}
\end{corollary}
\begin{proof}
Taking the product of squared differences in \eqref{m:hermitesplit}
gives the leading coefficient in \eqref{m:Dleading}.  Holomorphy in
$s$ removes any noninteger power from its remainder.  Substitution
in \eqref{m:ode}, using holomorphy of $B$, proves
\eqref{m:Tleading}.  The first equality in \eqref{m:charge} is the
finite-cluster charge formula.  For the second, integrate
\eqref{m:ode} divided by $\mathfrak D$; the integral of the holomorphic
coefficient $B$ is zero.
\end{proof}

The exterior contribution still changes $\mathfrak D$, $\mathfrak T$
and their zero locations through \eqref{m:ode}; its vanishing contour
integral is not permission to omit it from the evolution.  For $m>2$
the forcing vanishes at a full birth.  This removes the constant
forcing used by the two-root univalence argument.  A full birth has
charge $N$, while two real discriminant signs record at most its
parity.  Conjugate nonreal event pairs can contribute further even
charge without changing those signs.

\subsection{A cubic real-time monotonicity theorem}

Assume now that $a$ is real, $D$ is invariant under conjugation and
$H(\overline t,\overline z)=\overline{H(t,z)}$.  For $m=3$, depress
the cluster polynomial:
\[
 P(t,z)=(z-c(t))^3+p(t)(z-c(t))+q_0(t).
\]
The notation $q_0$ distinguishes the constant coefficient from $U_z/U$.
Direct expansion gives
\begin{equation}\label{m:cubicDT}
 \mathfrak D=-4p^3-27q_0^2,\qquad \mathfrak T=9p^2.
\end{equation}
For the second identity, $\mathfrak T$ is a translation-invariant
symmetric polynomial of degree four.  After centering, the fundamental
symmetric polynomials have degrees two and three, so it is a constant
multiple of $p^2$.  Substituting roots $(-1,0,1)$ gives the constant
$9$.

\begin{proposition}[At most one real cubic collision time]
\label{m:cubicmonotone}
On each real interval contained in $D$, a proper real-symmetric
three-root heat cluster has at most one collision time.  An
independently located collision at a strictly negative time in such an
interval excludes an additional collision at time zero in that
interval.
\end{proposition}
\begin{proof}
Choose an antiderivative of $B$ which is real on the real
interval, and put $\mu=\exp(-\int B)$.  Then $\mu>0$ there
and \eqref{m:ode}--\eqref{m:cubicDT} imply
\begin{equation}\label{m:cubicderivative}
                  (\mu\mathfrak D)'=72\mu p^2\geq0.
\end{equation}
If $p$ is not identically zero, its real zeros are isolated.  The
integral of the right side over every nontrivial subinterval is then
positive, so $\mu\mathfrak D$ is strictly increasing and has at most
one zero.  If $p$ is identically zero, $\mu\mathfrak D$ is constant
and nonzero, since $\mathfrak D$ cannot be identically zero.  There
are no collisions in that case.
\end{proof}

A real event in this cubic cluster can be a double, of charge one,
or a triple, of charge three.  Thus the proposition bounds the number
of real event locations, not their total charge.  It also gives no
bound on nonreal event times.  At complex time $p^2$ need not have
positive real part, even if $\mu$ lies in a fixed open half-plane.
For $m\geq4$, even the real inequality in
\eqref{m:cubicderivative} need not persist.

\subsection{Exact controls for the failure of the cap formula}

The following polynomial heat solutions pay the cluster and boundary
conditions explicitly.  They disprove the extension of the two-root
cap formula to arbitrary multiplicity, without challenging its
two-root hypothesis.  Their conclusions are algebraic; the exact
arithmetic replay in \cite{controls} is supplementary support.

\begin{proposition}[A cubic with three event times]
\label{m:cubiccontrol}
Let
\[
 H(t,z)=z^3-6tz+2,\qquad
 D=\{t:|\operatorname{Re}t|<1,\ |\operatorname{Im}t|<1\},
 \qquad r=4,\quad R=5,\quad a=0.
\]
This is a proper three-root heat cluster with $B=0$.  Its
total charge is three, although
$\bigl(\operatorname{sgn}\mathfrak D(1)
       -\operatorname{sgn}\mathfrak D(-1)\bigr)/2=1$.
\end{proposition}
\begin{proof}
One has $H_t=-6z=-H_{zz}$.  On the inner wall and closed time
rectangle,
\[
 |6tz-2|\leq24\sqrt2+2<64=|z^3|.
\]
Rouch\'e's theorem puts all three polynomial roots strictly in the
inner disk, hence also in the outer disk; both walls are nonzero.
Here $U=1$, and thus $B=0$.  Formula
\eqref{m:cubicDT} gives
\[
 \mathfrak D(t)=864t^3-108,
 \qquad \mathfrak T(t)=324t^2.
\]
The event times are the three simple roots of $t^3=1/8$:
\[
 \frac12,\qquad -\frac14+\frac{\sqrt3}{4}i,
 \qquad -\frac14-\frac{\sqrt3}{4}i.
\]
At each event the repeated root is $z=1/(2t)$, of modulus one;
$H_{zz}=6z\ne0$ there, so each event is an ordinary double.  All
three lie strictly within the product domain, and no event is on its
boundary.  Each has charge one.  Finally,
$\mathfrak D(-1)=-972$ and $\mathfrak D(1)=756$, proving the
claimed discrepancy.
\end{proof}

\begin{proposition}[A quadruple endpoint with equal positive cap signs]
\label{m:quarticcontrol}
On the same time rectangle, let
\[
 H(t,z)=z^4-12tz^2+12t^2,
 \qquad a=0,\quad r=5,\quad R=6.
\]
This proper four-root heat cluster has $B=0$, nonzero
time-boundary discriminant, equal positive real cap discriminants and
total charge six.  For every positive real time its spatial roots are
real and simple.
\end{proposition}
\begin{proof}
The heat equation follows by direct differentiation.  The inner-wall
estimate is
\[
 |12tz^2-12t^2|\leq300\sqrt2+24<625=|z^4|,
\]
so all four roots lie strictly in the inner disk throughout the
closed rectangle.  Again $U=1$ and $B=0$.  Set
$\alpha=6+2\sqrt6$, $\beta=6-2\sqrt6$.  The factorization
\[
 H(t,z)=(z^2-\alpha t)(z^2-\beta t)
\]
gives roots $\pm\sqrt{\alpha t}$ and $\pm\sqrt{\beta t}$.
The squared differences within the two opposite pairs contribute
$16\alpha\beta t^2$, and the four cross-pair squared differences
contribute $(\alpha-\beta)^4t^4$.  Therefore
\[
 \mathfrak D(t)=16\alpha\beta(\alpha-\beta)^4t^6
              =1{,}769{,}472\,t^6,
 \qquad
 \mathfrak T(t)=1{,}327{,}104\,t^5,
\]
where the last identity also follows from $\mathfrak D'=8\mathfrak T$.
The only event is $(0,0)$, of charge six.  Both cap discriminants are
positive, and no point of the time boundary is an event.  Since
$\alpha>\beta>0$, the four displayed roots are real and distinct for
$t>0$.  Notice also that $\mathfrak T$ changes sign on real time.
\end{proof}

More generally, the Hermite family $e^{-t\partial_z^2}z^m$ has
$\mathfrak D=d_mt^N$, $\mathfrak T=Nd_mt^{N-1}/8$ and
$B=0$; a sufficiently large spatial disk gives a proper
finite window.  Thus zero exterior phase alone is compatible with
every full-birth charge.  The controls above concern the specified
phase-and-cap inference.  They do not preclude other boundary
estimates or estimates using additional properties of the theta
source.

\subsection{Why centered variance cannot replace the discriminant}

Define
\[
 c=\frac1m\sum_i z_i,\qquad
 V=\sum_i(z_i-c)^2=\frac1m\sum_{i<j}(z_i-z_j)^2.
\]
The equality follows by expanding the pairwise sum.  This quantity is
holomorphic even when root labels are not.

\begin{proposition}[Centered second-moment evolution]\label{m:variance}
For a proper $m$-root heat cluster,
\begin{equation}\label{m:Vflow}
 V'=2m(m-1)+4\sum_i(z_i-c)q(t,z_i).
\end{equation}
At a full birth $(t_0,z_0)$,
\begin{equation}\label{m:Vbirth}
 V=2m(m-1)(t-t_0)+O((t-t_0)^2).
\end{equation}
\end{proposition}
\begin{proof}
Since $\sum_i(z_i-c)=0$, differentiating gives
$V'=2\sum_i(z_i-c)z_i'$.  Substitute \eqref{m:motion}.  For
each unordered internal pair the contribution is
\[
 4\left(\frac{z_i-c}{z_i-z_j}
       +\frac{z_j-c}{z_j-z_i}\right)=4.
\]
There are $N$ pairs, proving \eqref{m:Vflow} at simple fibers.  The
remaining symmetric sum is holomorphic, for it equals
\[
 \frac{1}{2\pi i}\oint_{|\zeta-a|=R}
       (\zeta-c)q(t,\zeta)\frac{P_z(t,\zeta)}{P(t,\zeta)}\,d\zeta.
\]
Continuation proves the identity everywhere.  At a full birth $V=0$
and each $z_i-c=0$, so evaluating \eqref{m:Vflow} gives
\eqref{m:Vbirth}.
\end{proof}

The real sign transition in \eqref{m:Vbirth} is compatible with
every Hermite birth.  Moreover, $V$ need not have the discriminant's
zero set: the cubic in Proposition~\ref{m:cubiccontrol} has
$V=12t$, hence $V(0)=0$ but $\mathfrak D(0)=-108\ne0$.
On an all-real fiber, $V=0$ does force full coalescence, but that
fiber condition cannot be assumed on an unknown negative cap.
Conversely, $V>0$ does not exclude a partial collision in a larger
real cluster.  Neither centered variance nor two discriminant signs
therefore replace the full multiplicity-sensitive event calculation.

\section{Supporting certificates and exact arithmetic checks}
\label{sec:support}

The analytic identities and counterexamples above do not depend on the
computations in this section. The numerical work serves two purposes:
to demonstrate that the complete theta source can satisfy the finite
boundary hypotheses, and to check implementations against exact models.
All reported numerical claims retain the scope and internal-review
status of their recorded calculations.

\subsection{A complete-theta application of the two-root theorem}

The interval calculation in \cite{note06,realcap} uses
\begin{equation}
 |\operatorname{Re}t|,|\operatorname{Im}t|\leq10^{-2},\qquad
 a=\frac{71}{2},\quad r=8,\quad R=10.
 \label{n:window}
\end{equation}
Both circular root counts are two. The enclosing data, rounded
outward for display, are
\begin{center}
\begin{tabular}{@{}lr@{}}
\toprule
Quantity & Certified bound\\
\midrule
Inner circular wall floor & $>8.36937\times10^{-7}$\\
Outer circular wall floor & $>7.57583\times10^{-7}$\\
Exterior coefficient modulus & $|A|<3.489078$\\
Phase product & $10^{-2}\sup|A|<0.034891$\\
Discriminant at each real time cap & $\Delta(\pm10^{-2})>144$\\
Complete source integrals & $4146$\\
\bottomrule
\end{tabular}
\end{center}
The two-root cap theorem therefore gives $Q=0$ on this window.
No RH assumption or tabulated zero location is an input.

The certificate covers entire arcs and the full complex-time rectangle,
not only sampled points. It uses 256 arcs on each circle, complete jets
of orders zero through fifteen at the time-zero anchors, and a
joint Taylor expansion of total degree seven. Complete moments pay the
order-eight remainder. All omitted source modes and exterior integration
tails are bounded. Zero-avoiding convex disks around the boundary image
arcs justify the polygonal winding counts. The coefficient estimate
retains the coupled numerator $H_z+(3/10)H$ before taking absolute
values. If $d_k$ and $n_k$ are the full-arc lower and upper bounds for
the denominator and this numerator, respectively, the calculation uses
\[
 |A|\leq\frac{4R}{(R-r)^2}\frac1{256}\sum_{k=0}^{255}\frac{n_k}{d_k}.
\]
At each real time cap, the certified signs at spatial points
$28,29,41,43$ are $+,-,-,+$. The intermediate value theorem and the
total count of two put one simple root in each of $(28,29)$ and
$(41,43)$. Their gap exceeds twelve, proving the displayed
discriminant bound. Thus the cap signs have an independent real-source
payment.

This is a separated-pair example. Its large cap gap does not address an
unknown shrinking gap near a hypothetical endpoint. The same boundary
and phase procedure also admits the exact quadratic collision
$H=(z-71/2)^2-2t$, returning charge one from its opposite cap signs.
It therefore does not silently assume that events are absent.

\subsection{Earlier direct comparisons}

The earlier finite experiments \cite{note01,note03} compared specified
boundary and direct joint-exclusion procedures using complete source
data. Their certified margins and acquisition counts were
\begin{center}
\begin{tabular}{@{}lccc@{}}
\toprule
Experiment & Boundary margin & Direct margin & Source integrals\\
\midrule
One-zero ellipsoid & $>7.495\times10^{-5}$ & $>7.993\times10^{-5}$ & $8$ versus $6$\\
Two-zero thin box & $>3.5683\times10^{-7}$ & $>3.9348\times10^{-7}$ & $1032$ each\\
\bottomrule
\end{tabular}
\end{center}
The thin box has $|\operatorname{Re}t|,|\operatorname{Im}t|\leq0.01$,
$28\leq\operatorname{Re}z\leq43$, and
$|\operatorname{Im}z|\leq0.01$. Both methods used the same source
acquisition and 2056 boundary vertices. The direct procedure had the
stronger certified margin in these declared implementations. This is
neither a universal complexity ordering nor evidence against other
boundary geometries. The analytic reference-transfer result in
Appendix~\ref{app:reference} explains why an already certified
holomorphic reference does not by itself confer an improved uniform
margin on the boundary method.

\subsection{Exact controls and reproducibility}

The standard-library rational calculation \cite{controls} checks the
Hermite discriminants for $2\leq m\leq8$, five finite clusters with
nonzero exterior forcing, the cubic expression for $T$, and quadratic
source response. It also verifies the explicit boundary-paid affine
homotopy example from the signed-response section. These checks agree
with the proofs, including the negative-time signs of the many-root
forcing. They are finite algebra checks, not a proof for all $m$.

The earlier piecewise affine boundary-degree controls \cite{note02}
return degrees $+1,-1,+1,0$ for, respectively, a complex-oriented
linear map, an orientation-reversing real map, an enclosed quadratic
heat event, and a quadratic heat field with its event outside the box.
The negative-degree example deliberately fails holomorphicity. It
checks the orientation convention rather than holomorphic positivity.

The repository records supply scripts, arithmetic versions, parameter
choices, source-tail and interpolation budgets, and small execution
receipts. The theta runs use \texttt{python-flint 0.9.0}; the exact
algebra checks require only Python's standard library. Source-array
hashes bind the recorded computations but do not replace their
analytic error estimates. No additional theta numerical region was
acquired for this manuscript.

\section{What remains to obtain endpoint exclusion}
\label{sec:open}

The geometric difficulty at an isolated hypothetical double has been
resolved in a pointwise quantitative sense: a positive radius, collar,
time width, pair of root counts, and phase margin follow from the
nonzero second jet and complete moments. The resulting window correctly
has charge one when a double is assumed. A new exclusion must therefore
use information beyond those universal local properties.

\paragraph{A sign for the unchanged theta source.}
The signed response and the canonical-tail mixture retain correlations
that absolute moment bounds discard. Their remaining task is to prove,
with a declared reference and the entire residual included, either a
strictly positive incoming discriminant or a strictly negative event
location. For a proper two-root window with a positive outgoing cap,
the first gives charge zero. Alternatively, an independently located
negative-time event uses the unique event allowed by the phase theorem
and excludes another event at zero. Neither required theta-specific
inequality has been established here. A strict improvement over another
enclosure can still leave zero inside the resulting interval.

\paragraph{Height coverage.}
The windows exist separately at every finite hypothetical double
height. Their sizes depend on the local curvature; the displayed
construction can shrink with height. The source sign, error terms,
tail cutoff and event-location margin must therefore be controlled with
the same dependence. A calculation on one fixed interval, or an
unquantified limiting statement about the head cutoff, cannot discharge
the universal quantifier over endpoint heights. No common negative
all-real time slice is being proposed.

\paragraph{Every multiplicity.}
The all-root identity $D'=8T+BD$ retains the whole cluster and its
exterior contribution. For $m>2$, the forcing $T$ is no longer constant.
The cubic control shows that nonreal time events survive even when
$B=0$; the quartic control has charge six despite equal positive real
cap signs. Applying the two-root theorem to an arbitrarily selected
pair at a higher birth is invalid, because the other birth roots
destroy the required isolating collar. The cubic real-event monotonicity
provides a possible conditional event-location argument, but no
complete-source location estimate is supplied. Higher multiplicities
still require an exclusion compatible with their full discriminant.

The finite-moment control constructed above has exactly its stated
preservation properties; it does not realize the actual theta moment
vector. Broader collision controls in the project have different
failures of the theta premises: Euler-factor controls change arithmetic
data; finite square-source and finite-prime controls lack the complete
normalized coefficient/reflection conjunction; reflected or smooth
positive-source controls lack the canonical square-source expansion;
the close positive endpoint controls in \cite{bw10control} fail ordinary
Dirichlet-series arithmetic. Their preservation lists must not be
combined. They rule out arguments using only their shared premises,
not estimates specific to the unchanged complete theta source.

In the task register, the hypothetical-double window is complete in
its stated parameterized scope (ECE-01). The distinguishing sign
(ECE-02), its all-height use (the signed part of ECE-04), and
higher-multiplicity exclusion (the remaining part of ECE-05) are open.
Endpoint exclusion and RH-conditional simplicity (ECE-03 and ECE-06)
remain contingent on those results. The manuscript is an analytic
consolidation, not an upgrade of those statuses.

\section*{Acknowledgements and research provenance}
\addcontentsline{toc}{section}{Acknowledgements and research provenance}
This manuscript was drafted for Edward Baker with substantial assistance
from OpenAI's Codex language-model system. The manuscript drafting and
internal review session recorded model \texttt{gpt-6-astra}, reasoning
effort \texttt{ultra}. The underlying 21 October 2026 notes record their
own configurations, including earlier \texttt{max} sessions; those
provenance distinctions are retained in the source records. Separate
agents of the same model assisted in drafting and checking sections.
These checks do not constitute independent specialist review or formal
verification. The author designation above identifies the intended
recipient of the draft and does not imply that he has approved every
claim. The investigation's notes, reviews, progress register and chat
archive preserve the detailed research history; no manuscript snapshot
series is created.

\appendix
\section{Holomorphic references and boundary interpolation}
\label{app:reference}

This appendix makes explicit two analytic ingredients of the earlier
finite comparisons \cite{note02}. They are separate from the signed
source estimates used in the endpoint program.

\begin{proposition}[Minimum norm for a zero-free square holomorphic map]
\label{a:minimum}
Let $\Omega\Subset\mathbb C^n$ be a bounded domain with piecewise smooth
boundary, and let $F$ be holomorphic near $\overline\Omega$, with
values in $\mathbb C^n$. Suppose
$b=\min_{\partial\Omega}\|F\|>0$. If
$\deg(F,\Omega,0)=0$, then $F$ has no zero in $\overline\Omega$ and
\[
 \min_{\overline\Omega}\|F\|=b.
\]
\end{proposition}
\begin{proof}
For $\|v\|<b$, the fiber $F^{-1}(v)\cap\Omega$ is a compact analytic
subset of the open Euclidean domain: it cannot approach the boundary.
Such a subset is finite. Indeed, on a compact irreducible analytic
component the maximum principle makes every holomorphic coordinate
function constant, excluding positive dimension; compactness then
makes the discrete set finite. At each isolated holomorphic preimage
the local degree is a positive integer, the analytic multiplicity.
The degree at $v$ is their sum. It is constant for $\|v\|<b$, since
that ball misses $F(\partial\Omega)$. Degree zero at $0$ therefore
implies that every such fiber is empty. This proves the claimed
interior lower bound, and compactness of the boundary gives equality.
\end{proof}

\begin{corollary}[Transfer of the same certified margin]
\label{a:transfer}
Under the hypotheses of Proposition~\ref{a:minimum}, let $G$ be another
square holomorphic map near $\overline\Omega$ and suppose
$\sup_{\partial\Omega}\|G-F\|\leq e<b$. Then
\[
 \|G\|\geq b-e\quad\hbox{on }\overline\Omega,
 \qquad \deg(G,\Omega,0)=0.
\]
\end{corollary}
\begin{proof}
Apply the scalar maximum-modulus principle to every unit linear
functional of $G-F$. Taking the supremum over those functionals gives
$\sup_{\overline\Omega}\|G-F\|\leq e$. The norm inequality follows
from Proposition~\ref{a:minimum}; the straight boundary homotopy
$F+s(G-F)$ is nonzero and preserves degree.
\end{proof}

The same supplied reference-degree proof, boundary floor and error
certificate thus yield a direct interior joint margin without additional
source acquisition. This is a comparison of those data, not a lower
bound for the cost of every boundary algorithm. Computing the reference
degree has a cost shared by both uses.

The margin is sharp even for a heat jet requiring both components.
On $|\operatorname{Re}t|,|\operatorname{Im}t|<1/10$ and
$|\operatorname{Re}z|,|\operatorname{Im}z|<3/2$, set
\[
 F=(z^2-1-2t,2z),\qquad
 G=(z^2-1-2t+\eta,2z),\quad 0<\eta<4/5.
\]
Writing $s=|z|^2$ and $k=1+2t$ gives
\[
 \|F\|^2\geq(s-|k|)^2+4s
 =|k|^2+(4-2|k|)s+s^2\geq16/25,
\]
because $\operatorname{Re}k\geq4/5$ and $|k|<2$.
Equality holds at $t=-1/10,z=0$.
The identical argument with $k-\eta$ gives exact floor $4/5-\eta$
for $G$, while $\|G-F\|=\eta$. At time zero, the first component
of $F$ vanishes at $z=\pm1$, and its second component at $z=0$.
Neither component by itself supplies the joint floor.

\subsection{A complete interpolation payment}

Suppose a convex boundary simplex is approximated by affine
interpolation of the jets $H_j=\partial_z^jH$, $j=0,1$.
Let $(t_i,z_i)$ be its vertices, $\lambda_i$ barycentric weights,
and $(t,z)=\sum_i\lambda_i(t_i,z_i)$. Put
\[
 V_t=\sum_i\lambda_i|t_i-t|^2,\qquad
 V_z=\sum_i\lambda_i|z_i-z|^2.
\]
If complete moments $M_k(T,R)$ bound all segment derivatives, Taylor's
formula and cancellation of the linear barycentric term give
\begin{equation}
 \left|\sum_i\lambda_iH_j(t_i,z_i)-H_j(t,z)\right|
 \leq I_j:=\frac12\left(M_{j+4}V_t+
 2M_{j+3}\sqrt{V_tV_z}+M_{j+2}V_z\right).
 \label{a:interpolation}
\end{equation}
To check this, the second directional derivative has terms
$H_{j,tt}=H_{j+4}$, $H_{j,tz}=-H_{j+3}$ and
$H_{j,zz}=H_{j+2}$. Integrating the Taylor remainder contributes
$1/2$, and Cauchy--Schwarz bounds the mixed barycentric sum.
For a real coordinate with range length $h$, its variance is at most
$h^2/4$: this follows by taking expectations in
$(X-a)(b-X)\geq0$ and maximizing
$(\mathbb EX-a)(b-\mathbb EX)$. Summing over the two time and
two space coordinates bounds $V_t,V_z$ explicitly.

If the vertex errors in $H_j$ are at most $\delta_j$, convexity keeps
the interpolated acquisition errors at most $\delta_j$. Thus the
full Euclidean jet error on the simplex is at most
\begin{equation}
 e_\sigma=\sqrt{(\delta_0+I_0)^2+(\delta_1+I_1)^2}.
 \label{a:fullerror}
\end{equation}
A certified lower bound $b_\sigma$ for the norm on the entire affine
image, with $e_\sigma<b_\sigma$, gives a nonzero straight boundary
homotopy. Vertex norms alone are insufficient for $b_\sigma$; for
example, enclosing the image convex hull in a coordinate box permits
a lower bound by that box's distance from zero.

For completeness, a finite piecewise affine degree computation can
cone a coherently triangulated box boundary to its center, assigning
the artificial target value zero to the cone vertex. This is an
extension of the boundary map, not an interior approximation to $G$.
For each four-simplex form the domain and target matrices
$X=[x_0\ x_1\ x_2\ x_3]$ and
$Y=[y_0\ y_1\ y_2\ y_3]$ relative to the cone center. A sufficiently
small regular target $v$ contributes $\operatorname{sgn}(\det Y/\det X)$
exactly when $Y^{-1}v$ has positive coordinates with sum less than one.
Choose $\|v\|<\min_\sigma b_\sigma$ and verify that it is on no
simplex image face and has no preimage in a singular simplex. The sum
is the boundary degree, by additivity and homotopy invariance. Rational
linear algebra can verify these finite conditions exactly. The
artificial center uses no evaluation of the actual heat field. Finally,
holomorphic positivity is what makes zero degree exclude a zero of
$(H,H_z)$; that last conclusion would fail for a general real map.

\begin{thebibliography}{99}
\addcontentsline{toc}{section}{References and supporting research records}

\bibitem{polymath}
D. H. J. Polymath,
\emph{Effective approximation of heat flow evolution of the Riemann
$\xi$ function, and a new upper bound for the de Bruijn--Newman constant},
arXiv:1904.12438v2 (2019), especially Sections 1 and 3.
\href{https://arxiv.org/abs/1904.12438v2}{Public preprint}.

\bibitem{rodgers-tao}
B. Rodgers and T. Tao,
\emph{The de Bruijn--Newman constant is non-negative},
arXiv:1801.05914v5 (2021).
\href{https://arxiv.org/abs/1801.05914v5}{Public preprint}.

\bibitem{note01}
BW12 Research Note 01,
\emph{Bounded theta charge and direct comparison}, 21 October 2026.
\href{../notes/01_BOUNDED_THETA_CHARGE_AND_DIRECT_COMPARISON_20261021.md}{Analytic error bounds and pilot certificate}.

\bibitem{note02}
BW12 Research Note 02,
\emph{Holomorphic reference transfer and a boundary degree prototype},
21 October 2026.
\href{../notes/02_HOLOMORPHIC_REFERENCE_TRANSFER_AND_BOUNDARY_DEGREE_20261021.md}{Proof and exact boundary controls}.

\bibitem{note03}
BW12 Research Note 03,
\emph{Two-zero theta boundary and direct comparison}, 21 October 2026.
\href{../notes/03_TWO_ZERO_THETA_BOUNDARY_AND_DIRECT_COMPARISON_20261021.md}{Comparison and complete-source payments}.

\bibitem{note04}
BW12 Research Note 04,
\emph{Analytical scope and potential of relative charge}, 21 October 2026.
\href{../notes/04_ANALYTICAL_SCOPE_AND_POTENTIAL_OF_RELATIVE_CHARGE_20261021.md}{Cluster and charge development}.

\bibitem{note05}
BW12 Research Note 05,
\emph{Two-root discriminant flow and real cap charge}, 21 October 2026.
\href{../notes/05_TWO_ROOT_DISCRIMINANT_FLOW_AND_REAL_CAP_CHARGE_20261021.md}{Two-root theorem and source interface}.

\bibitem{note06}
BW12 Research Note 06,
\emph{Complete theta walls and a real cap charge application},
21 October 2026.
\href{../notes/06_COMPLETE_THETA_WALLS_AND_REAL_CAP_APPLICATION_20261021.md}{Finite interval-certified application}.

\bibitem{note07}
BW12 Research Note 07,
\emph{Quantitative endpoint windows and height coverage}, 21 October 2026.
\href{../notes/07_QUANTITATIVE_ENDPOINT_WINDOWS_AND_HEIGHT_COVERAGE_20261021.md}{Parameterized window proof}.

\bibitem{note08}
BW12 Research Note 08,
\emph{Signed discriminant response and event displacement},
including Appendix A, 21 October 2026.
\href{../notes/08_SIGNED_DISCRIMINANT_RESPONSE_AND_EVENT_DISPLACEMENT_20261021.md}{Response, controls and canonical-tail refinement}.

\bibitem{note09}
BW12 Research Note 09,
\emph{Many-root discriminant flow and cap limits}, 21 October 2026.
\href{../notes/09_MANY_ROOT_DISCRIMINANT_FLOW_AND_CAP_LIMITS_20261021.md}{All-root identities and exact controls}.

\bibitem{review08}
BW12 Review 08,
\emph{Endpoint charge opening-round review}, 21 October 2026.
\href{../reviews/08_ENDPOINT_CHARGE_OPENING_ROUND_REVIEW_20261021.md}{Same-model internal review record}.

\bibitem{tail04}
BW14 Research Note 04,
\emph{Canonical tail Laplace constraint and coupled enclosure},
21 October 2026.
\href{../../14_theta_source_transform_constraints/notes/04_CANONICAL_TAIL_LAPLACE_CONSTRAINT_AND_COUPLED_ENCLOSURE_20261021.md}{Explicit positive tail measure}.

\bibitem{tail05}
BW14 Research Note 05,
\emph{Finite reflection jets and head-tail stability}, 21 October 2026.
\href{../../14_theta_source_transform_constraints/notes/05_FINITE_REFLECTION_JETS_AND_HEAD_TAIL_STABILITY_20261021.md}{Reference-change and residual qualifications}.

\bibitem{bw10}
\emph{RH-Conditional Approximation Estimates and the Remaining
Simplicity Problem}, working manuscript drafted for Edward Baker.
\href{../../10_rh_implies_simple_zeros/manuscript/main.tex}{Companion BW10 manuscript}.

\bibitem{assessment07}
BW12 Review 07,
\emph{Assessment of the connection to BW10 and the proof direction},
21 October 2026.
\href{../reviews/07_BW10_INTEGRATION_AND_PROOF_DIRECTION_ASSESSMENT_20261021.md}{Weighted endpoint-charge identity and scope}.

\bibitem{bw10control}
BW10 Research Note 05,
\emph{Close positive sources and normalized double zeros},
20 October 2026.
\href{../../10_rh_implies_simple_zeros/notes/05_CLOSE_POSITIVE_SOURCES_AND_NORMALIZED_DOUBLE_ZEROS_20261020.md}{Conditional endpoint control and its arithmetic limitation}.

\bibitem{realcap}
BW12 numerical records,
\href{../numerics/REAL_CAP_THETA_APPLICATION_20261021.json}{Real-cap theta application receipt}
and \href{../numerics/certify_real_cap_application.py}{replay verifier},
21 October 2026. Dependencies and execution instructions are in the
\href{../numerics/README.md}{numerics guide}.

\bibitem{controls}
BW12 numerical records,
\href{../numerics/ECE_EXACT_CONTROLS_20261021.json}{Endpoint-program exact algebra receipt}
and \href{../numerics/ece_exact_controls.py}{exact rational replay},
21 October 2026.

\end{thebibliography}

\end{document}
